\chapter{傅里叶变换和拉普拉斯变换}

\begin{flushright}
\textit{“往来不穷谓之通。”——《周易·系辞上》}
\footnote{积分变换把函数送进另一个域,反变换再送回来.两边合在一起,微分变成乘法,卷积变成乘积,初值问题变成代数.}
\end{flushright}

这一章我们将会学习非常有用的\term{傅里叶级数}{fuliyejishu}
(Fourier Series)和\term{傅里叶变换}{fuliyebianhuan}(Fourier Transform)~\cite{fourier,bracewell,lighthill}.
傅里叶级数和傅里叶变换在数论、组合数学、信号处理、概率论、统计学、密码学、声学、光学等领域都有着广泛的应用.
拉普拉斯(Laplace)变换在光学等工程技术与科学领域中有着广泛的应用
% \subsection{傅里叶变换}
\label{subsec:fourier_transform}
前面我们讨论了周期性体系的傅里叶级数展开.现在我们希望研究非周期性函数的傅里叶展开问题.
假设$f(x)$是定义在$-\infty < x < \infty$的实函数,我们依旧对该函数进行傅里叶级数展开,
该级数必须理解为周期$2\ell$扩展至$\infty$的极限情况.
由于三角函数族的变量为
\[
 \frac{n\pi x}{\ell}    
\]
引入非连续变量
\[\omega_n = \frac{n\pi} {\ell}, \quad k = 0,1,2,\cdots
    \]
可知$\omega_n = n \omega_0$, 其中$\omega_0 = \frac{\pi}{\ell}$为基础频率(Fundamental Frequency),
\[\Delta \omega_n = \omega_n - \omega_{n-1} = \omega_0 . \] 
有
\begin{equation}
    f(x) = \lim_{\ell\to \infty} \sum_{n=-\infty}^{\infty} \left[
        c_n e^{\imath \omega_n x}
        % \frac{a_0}{2} + \sum_{n=1}^{\infty} \left( a_n \cos \omega_n x + b_n \sin \omega_n x \right)
    \right]
\end{equation}
其中系数
\begin{equation}
    c_n = \frac{1}{2\ell} \int_{-\ell}^{\ell} f(x) e^{-\imath \omega_n x} dx 
\end{equation}
取$\ell\to \infty$的极限,则
\[
  \sum_n \to  \frac{ \ell}{\pi}  \int d\omega 
\]
傅里叶系数记为$F(\omega)$,
有
\begin{equation}
    f(x) = \int_{-\infty}^{\infty} d \omega F(\omega) e^{\imath \omega x}, % \quad F (\omega) \equiv 2\ell c_n
\end{equation}
上式的积分称为\textbf{傅里叶积分}(Fourier integral),是$f(x)$的\textbf{傅里叶变换}(Fourier transform).
其中
\begin{equation}
    F (\omega) \equiv \lim_{\ell\to \infty} \frac{ 2\ell c_n}{2\pi}  = \frac{1}{2\pi} \int_{-\infty}^{\infty} e^{-\imath \omega x} f(x) dx,
\end{equation}
这被称为\textbf{逆傅里叶变换}(inverse Fourier transform). 这里积分前的系数并不要紧,更关键的是积分关于$\omega$的形式.

傅里叶积分存在的条件由\textbf{傅里叶积分定理}判定: 
若函数 $f(x)$ 在区间 $(-\infty, \infty)$ 上满足条件: $(1) f(x)$ 在 任一有限区间上满足狄里希利条件; (2) $f(x)$ 在 $(-\infty, \infty)$ 
上绝对可积 (即 $\int_{-\infty}^{\infty}|f(x)| dx$ 收敛), 则 $f(x)$ 可表示成傅里叶积分, 且
$$
\textrm{傅里叶积分值} = \lim_{\epsilon\to 0} \half \left[ f(x+\epsilon) + f(x-\epsilon) . \right]
$$
% $F (\omega) \equiv 2\ell c_n$.

% \begin{align}
%     a_n =  \lim_{\ell\to \infty} \frac{1}{\ell} \int_{-\ell}^{\ell} f(x) \cos {  \omega_n x  } dx
%     \\
%     b_n = \frac{1}{\ell} \int_{-\ell}^{\ell} f(x) \sin {   \omega_n x } dx
% \end{align}

正余弦形式利用 $c_n$ 与 $a_n$、$b_n$ 的关系得到.
上一节已有 ($n>0$)
\[
  c_n=\half(a_n-\imath b_n),
  \qquad
  c_{-n}=\half(a_n+\imath b_n),
\]
亦即 $a_n=c_n+c_{-n}$, $b_n=\imath(c_n-c_{-n})$.
把 $\pm n$ 两项合并:
\[
  c_n e^{\imath\omega_n x}+c_{-n}e^{-\imath\omega_n x}
  =a_n\cos(\omega_n x)+b_n\sin(\omega_n x).
\]
又 $F(\omega)=\lim(2\ell c_n)/(2\pi)=\lim c_n/\Delta\omega$, 相应定义
\[
  A(\omega)=\lim\frac{a_n}{\Delta\omega}=F(\omega)+F(-\omega),
  \quad
  B(\omega)=\lim\frac{b_n}{\Delta\omega}=\imath\bigl[F(\omega)-F(-\omega)\bigr],
\]
或等价地 $F(\omega)=\half\bigl[A(\omega)-\imath B(\omega)\bigr]$.
由 $a_n=(1/\ell)\int f\cos$ 与 $1/\ell=\Delta\omega/\pi$, 取 $\ell\to\infty$ 得
\begin{equation}
    \left\{\begin{array}{l}
    A(\omega)=\frac{1}{\pi} \int_{-\infty}^{\infty} f(t) \cos \omega t\, d t, \\
    B(\omega)=\frac{1}{\pi} \int_{-\infty}^{\infty} f(t) \sin \omega t\, d t .
    \end{array}\right.
\end{equation}
级数中 $\pm n$ 合并后再取极限, 得到正余弦形式
\begin{equation}
    f(x) =\int_{0}^{\infty} A(\omega) \cos {\omega x} d\omega + \int_{0}^{\infty} B(\omega) \sin {\omega x} d\omega.
\end{equation}

上式还可以写成

\begin{equation}
    f(x)=\int_0^{\infty} C(\omega) \cos [\omega x-\varphi(\omega)] d \omega,
\end{equation}

其中
$$
\begin{gathered}
C(\omega)=\sqrt{[A(\omega)]^2+[B(\omega)]^2} \\
\varphi(\omega)=\arctan [B(\omega) / A(\omega)] .
\end{gathered}
$$
$C(\omega)$ 称为 $f(x)$ 的振幅谱, $\varphi(\omega)$ 称为 $f(x)$ 的相位谱.如果$f(x)$是奇函数或偶函数,
对应的级数为\textbf{傅里叶余弦积分}或\textbf{傅里叶正弦积分}.

复数形式的傅里叶积分可以写成对称的形式,
\begin{equation}
    \begin{aligned}
    f(x) & =\frac{1}{\sqrt{2 \pi}} \int_{-\infty}^{\infty} F(\omega)e^{\imath \omega x} d \omega, \\
    F(\omega) & =\frac{1}{\sqrt{2 \pi}} \int_{-\infty}^{\infty} f(x)\left[e^{\imath \omega x}\right]^* d x .
    \end{aligned}
\end{equation}
并常用符号简写为
$$
F(\omega)=\mathcal{F}[f(x)], \quad f(x)=\mathcal{F}^{-1}[F(\omega)] .
$$
$f(x)$ 和 $F(\omega)$ 分别称为傅里叶变换的\textbf{原函数}和\textbf{像函数}.

\begin{example}
求以下函数的复数傅里叶变换.
\begin{enumerate}
    \item $f(x) = e^{-\alpha |x|}, \alpha > 0$;
    \item $f(x) = \delta (x)$;
    \item $f(x) = e^{-\alpha x^2}, \alpha > 0$.
\end{enumerate}
\end{example}
\begin{solution}
\begin{enumerate}
    \item 对$x$分段进行处理,利用对称形式的傅里叶变换有
    \begin{align}
        F(\omega) &=
        \sqrt{\frac{1}{2 \pi}} \int_{-\infty}^0 e^{\alpha x+i \omega t} dx \nonumber
        % \\ &
        +
        \sqrt{\frac{1}{2 \pi}} \int_0^{\infty} e^{-\alpha x+i \omega t} dx  \nonumber
        \\
        &= \sqrt{\frac{1}{2 \pi}} \left[\frac{1}{\alpha + \imath \omega } + \frac{1}{\alpha - \imath \omega } \right] \nonumber
        \\
        &= \sqrt{\frac{1}{2 \pi}} \frac{2\alpha}{\alpha^2 + \omega^2}, \nonumber
    \end{align}
    越大的$\alpha$则$f(x)$越窄,集中在$x=0$附近,也就是越局域,而其傅里叶变换则越弥散.当$f(x)$是偶函数时,傅里叶变换的像函数为实函数.

    \item 代入
    \begin{align}
        F(\omega) &=
        \sqrt{\frac{1}{2 \pi}} \int_{-\infty}^{\infty} \delta(x) e^{+\imath \omega x} d x = \sqrt{\frac{1}{2 \pi}} ,
    \end{align}
    $f(x)$无限局域对应的傅里叶变换为无限弥散.

    \item 类似的,  
    \begin{align}
        F(\omega) &=\frac{1}{\sqrt{2 \pi}} \int_{-\infty}^{\infty} e^{-\alpha x^2} e^{\imath \omega x} d x \nonumber
        % \\ &
        % \\
        % & 
        = \frac{ e^{-\omega^2 / 4\alpha }}{\sqrt{2 \pi}} \int_{-\infty}^{\infty} e^{-\alpha  \left( x - \frac{\imath \omega}{2\alpha} \right)^2}  d x
        \nonumber
        % \\ & 
        \\
        & = \frac{ e^{-\omega^2 / 4\alpha }}{\sqrt{2 \pi}} \int_{-\infty - \imath \omega / 2\alpha}^{\infty - \imath \omega / 2\alpha} e^{-\alpha t^2} dt
        \nonumber
        % \\ &
        \\
        & = \frac{ e^{-\omega^2 / 4\alpha }}{\sqrt{2 \pi}} \sqrt{\frac{\pi}{\alpha}} 
        \nonumber
        % \\ &
        \\
            & = \frac{1}{\sqrt{2\alpha}} e^{- \frac{\omega^2}{4\alpha}} \nonumber .
        % \\ &
    \end{align}
    上式用了变量代换$t = x - \imath \omega /2 \alpha$,高斯函数的像函数还是高斯函数.
\end{enumerate}
\end{solution}


\subsection{傅里叶变换基本性质}
\idx{卷积}{juanji}
傅里叶变换满足以下基本性质.假定$f(x)$的傅里叶变换存在,记为$\mathcal{F}[f(x)] = F(\omega)$.

\begin{enumerate}
    \item 导数定理
    \begin{equation}
        \mathcal{F} [f'(x)] = \imath \omega F(\omega)
    \end{equation}

    \item 积分定理
    \begin{equation}
        \mathcal{F} [ \int^{x} f(x) dx ] = \frac{1}{\imath \omega} F(\omega)
    \end{equation}

    \item 相似性定理
    \begin{equation}
        \mathcal{F} [ f(ax) ] = \frac{1}{|a|} F\Bigl(\frac{\omega}{a}\Bigr),
        \qquad a\neq 0.
    \end{equation}
    绝对值来自雅可比. $a<0$ 时坐标反向,写成 $1/a$ 会丢一个负号.
    \item 延迟定理
    \begin{equation}
        \mathcal{F} [ f(x - x_0 ) ] = e^{-\imath \omega x_0} F(\omega),
    \end{equation}
    \item 位移定理
    \begin{equation}
        \mathcal{F} [ e^{\imath \omega_0 x} f(x) ] = F(\omega - \omega_0),
    \end{equation}
    \item 卷积定理. 卷积定义为
    \begin{equation}
        (f_1 * f_2)(x) = \int_{-\infty}^{\infty} f_1(\alpha)\, f_2(x-\alpha)\, d\alpha .
    \end{equation}
    注意第二个因子的自变量是 $x-\alpha$,不是 $\alpha-x$.在本节的对称约定下,
    \begin{equation}
        \mathcal{F}[f_1 * f_2] = \sqrt{2\pi}\, F_1(\omega)\, F_2(\omega).
    \end{equation}
    因子 $\sqrt{2\pi}$ 来自归一化.若把 $2\pi$ 全部放在逆变换上,这个常数会换成别的数,乘积结构不变.
\end{enumerate}

导数定理可以由分部积分直接看出.边界项在无穷远处消失时,
\begin{equation}
    \int_{-\infty}^{\infty} f'(x)\, e^{-\imath \omega x}\,\mathrm{d}x
    = \imath\omega \int_{-\infty}^{\infty} f(x)\, e^{-\imath \omega x}\,\mathrm{d}x,
\end{equation}
两边再除以 $\sqrt{2\pi}$ 就是 $\mathcal{F}[f']=\imath\omega F$.延迟定理是换元 $u=x-x_0$:指数多出一个 $e^{-\imath\omega x_0}$,积分变回 $F(\omega)$.

\begin{example}
求宽度为 $2a$ 的一维狭缝 $f(x)=1$（$|x|\le a$）的傅里叶变换,并由此说明矩孔.
\end{example}
\begin{solution}
\begin{equation}
    F(\omega)
    = \frac{1}{\sqrt{2\pi}}\int_{-a}^{a} e^{-\imath\omega x}\,\mathrm{d}x
    = \sqrt{\frac{2}{\pi}}\, a\, \frac{\sin(a\omega)}{a\omega}.
\end{equation}
$\omega=0$ 时商式的极限是 $\sqrt{2/\pi}\,a$.缝越宽,中央峰越窄.二维矩孔是两个这样的因子相乘.
\end{solution}

\subsection{高维傅里叶变换}

二维连续函数 $f(x, y)$ 的傅里叶变换定义如下:
设 $f(x, y)$ 是两个独立变量 $x, y$ 的函数, 且在 $\pm \infty$ 上绝对可积, 则定义积分
$$
F\left(k_1, k_2\right)=\frac{1}{2 \pi} \int_{-\infty}^{\infty} \int_{-\infty}^{\infty} f(x, y) e^{-\imath\left(k_1 x+k_2 y\right)} d x d y
$$
为二维连续函数 $f(x, y)$ 的傅里叶变换,并定义
$$
f(x, y)=\int_{-\infty}^{\infty} \int_{-\infty}^{\infty} F\left(k_1, k_2\right) e^{\imath\left(k_1 x+k_2 y\right)} d k_1 d k_2
$$
为 $F\left(k_1, k_2\right)$ 的逆变换.
这是把 $2\pi$ 全部放在正变换上的写法.与一维例子相同的对称约定是两侧各有一个 $1/(2\pi)$:
\begin{equation}
    F(k_1,k_2)=\frac{1}{2\pi}\iint f(x,y)\,e^{-\imath(k_1 x+k_2 y)}\,\mathrm{d}x\,\mathrm{d}y,
\end{equation}
\begin{equation}
    f(x,y)=\frac{1}{2\pi}\iint F(k_1,k_2)\,e^{\imath(k_1 x+k_2 y)}\,\mathrm{d}k_1\,\mathrm{d}k_2.
\end{equation}
两种约定不能在同一次计算里混用.下面的矩孔按对称约定计算.

\begin{example}
    求函数 $f(x, y)= \begin{cases}A, & |x| \leq X,|y| \leq Y \\ 0, & |x|>X,|y|>Y\end{cases}$ 的傅里叶变换 (矩孔费琅和夫衍射).
\end{example}
\begin{solution}
    由傅里叶变换关系
$$
F\left(k_1, k_2\right)=\frac{1}{2 \pi} \int_{-\infty}^{\infty} \int_{-\infty}^{\infty} f(x, y) e^{-i\left(k_1 x+k_2 y\right)} d x d y
$$
有
$$
\begin{aligned}
F\left(k_1, k_2\right) & =\frac{A}{2 \pi} \int_{-X}^X e^{-i k_1 x} d x \int_{-Y}^Y e^{-i k_2 y} d y \\
& =\left.\left.\frac{A}{2 \pi} \frac{1}{-i k_1} e^{-i k_1 x}\right|_{-X} ^X \frac{1}{-i k_2} e^{-i k_2 y}\right|_{-Y} ^Y \\
& =\frac{2A}{\pi} \frac{1}{i 2 k_1}\left(e^{i k_1 X}-e^{-i k_1 X}\right) \frac{1}{i 2 k_2}\left(e^{i k_2 Y}-e^{-i k_2 Y}\right) \\
& =\frac{2A X Y}{\pi} \frac{\sin \left(k_1 X\right)}{k_1 X} \frac{\sin \left(k_2 Y\right)}{k_2 Y}
\end{aligned}
$$
\end{solution}

物理上,$f(x,y)$ 是孔径面上的复振幅透过率:矩形开口内均匀照明取常值 $A$,挡光处为零.
费琅和夫(远场)衍射振幅正是孔径函数的傅里叶变换 $F(k_1,k_2)$,观察屏上的光强正比于 $|F|^2$;
$k_1,k_2$ 对应衍射角方向.
轴上 $F(0,0)=2AXY/\pi$,与孔径面积成正比.
$k_1 X=n\pi$ 或 $k_2 Y=m\pi$($n,m=\pm 1,\pm 2,\ldots$) 处振幅为零,屏上出现互相垂直的暗纹.
孔越窄,中央峰在对应方向越宽,与一维狭缝同一条宽窄互换.
$F$ 又是两个一维狭缝因子的乘积,故矩孔图样是互相垂直的两条单缝衍射结构.

三维时对称约定是两侧各有 $(2\pi)^{-3/2}$:
\begin{equation}
    F(\vec{k})=\frac{1}{(2 \pi)^{3/2}} \iiint_{-\infty}^{\infty} f(\vec{r}) e^{-\imath  \vec{k} \cdot \vec{r}} d^3 \vec{r},
\end{equation}
\begin{equation}
     f(\vec{r})=\frac{1}{(2 \pi)^{3/2}}\iiint_{-\infty}^{\infty} F(\vec{k}) e^{\imath \vec{k} \cdot \vec{r}} d^3 \vec{k}.
\end{equation}
把 $(2\pi)^{-3}$ 只放在一侧是另一种等价约定,不能和上式混用.指数可分离时,积分拆成三个一维积分.

若 $f(\vec{r})=f(r)$ 球对称,则 $F$ 只依赖 $k=|\vec{k}|$.
取极轴沿 $\vec{k}$,角积分给出 $4\pi\sin(kr)/(kr)$,对称约定下
\begin{equation}
    F(k)
    =\sqrt{\frac{2}{\pi}}\,\frac{1}{k}
      \int_0^{\infty} r\,f(r)\sin(kr)\,\mathrm{d}r.
\end{equation}
三维球对称傅里叶变换化成一维正弦变换.

\begin{example}
求汤川势 $f(r)=e^{-\mu r}/r$($\mu>0$) 的三维傅里叶变换,并说明它在散射中的用法.
\end{example}
\begin{solution}
用 $\int_0^{\infty}e^{-\mu r}\sin(kr)\,\mathrm{d}r=k/(k^2+\mu^2)$,得
\begin{equation}
    F(k)=\sqrt{\frac{2}{\pi}}\,\frac{1}{k^2+\mu^2}.
\end{equation}
散射计算里动量空间矩阵元常不带 $(2\pi)^{-3/2}$,直接写
\begin{equation}
    \int\frac{e^{-\mu r}}{r}\,e^{-\imath\vec{q}\cdot\vec{r}}\,\mathrm{d}^3\vec{r}
    =\frac{4\pi}{q^2+\mu^2},
    \qquad \vec{q}=\vec{k}'-\vec{k}.
\end{equation}
这是玻恩近似下一阶散射振幅的势因子.$\mu\to 0$ 给出库仑势的 $4\pi/q^2$,对应卢瑟福公式的动量依赖.
\end{solution}
由于它的像原函数$f(x)$要求的条件比傅里叶变换的条件要弱,因此在某些问题上,它比傅里叶变换的适用面要广.
本章首先从傅里叶变换的定义出发,导出拉普拉斯变换的定义,并研究它的一些基本性质,然后给出其逆变换的积分表达式--复反演积分公式,
并得出像原函数的求法,最后介绍拉普拉斯变换的应用．

\section{傅里叶级数}
\label{sec:fourier_series}
\subsection{正交函数组}
\idx{正交函数组}{zhengjiaohanshuzu}
\begin{definition}
  \label{def:orthogonal_function}
  设$f(x), g(x) $为定义在区间$[a,b]$上的函数, 积分
  \begin{equation}
    (f, g) \equiv \int_a^b f(x) g(x) dx
  \end{equation}
叫做二函数$f(x)$和$g(x)$的内积$(f,g)$或$(fg)$.
\end{definition}
内积满足施瓦茨不等式(Schwarz inequality)
\begin{equation}
  |(f,g)|^2 \leq (f,f)(g,g)
\end{equation}
当且仅当$f(x)$和$g(x)$线性相关时,等号成立.
如果$(f,g) = 0$,则称$f(x)$和$g(x)$是\textbf{正交}的. 函数$f(x)$和它自己的内积称为\textbf{范数}或\textbf{模},记作
$N f$,
\begin{equation}
  N f \equiv (f,f) = \int  f(x)^2 dx
\end{equation}
范数为1的函数称为\textbf{单位函数}或\textbf{归一化函数}. 一组归一化函数$\varphi_1(x), \varphi_2(x), \cdots, \varphi_n(x)$,其中任意二者皆正交，
称\textbf{正交归一组}, 表示为
\begin{equation}
  (\varphi_i, \varphi_j) = \delta_{ij} .
\end{equation}
对于复变函数， 内积定义为
\begin{equation}
  ( f, g ) = \int  f^*(x) g(x) dx
\end{equation}
其中$\bar f$为$f$的复共轭.  若$N f = \int | f|^2 dx = 1$, 称函数$f(x)$是归一化的. 

若对于所有$x$都成立的常系数齐次线性关系
\begin{equation}
  \sum_{i=1}^n c_i \varphi_i(x) = 0,
\end{equation}
其中系数不全为零, 我们称这$n$个函数$f_1, f_2, \cdots, f_n$是\textbf{线性相关}的. 否则称为\textbf{线性无关}的. 值得注意的是
正交组中的各函数永远是线性无关的. 因为, 若恒等式
\begin{equation}
  \sum_{i=1}^n c_i \varphi_i(x) = 0,
\end{equation}
成立，我们可以左乘$\varphi_j(x)$并积分可得$c_j = 0$.

我们可以类比于线性代数中向量的施密特正交化的方法对函数组进行正交化. 设$f_1(x), f_2(x), \cdots, f_n(x)$是线性无关的函数组, 我们可以通过
\begin{equation}
  \varphi_1(x) = f_1(x), \quad \varphi_2(x) = f_2(x) - \frac{(f_2, \varphi_1)}{(\varphi_1, \varphi_1)} \varphi_1(x), \cdots
  \varphi_n(x) = f_n(x) - \sum_{j=1}^{n-1} \frac{(f_n, \varphi_j)}{(\varphi_j, \varphi_j)} \varphi_j(x)
\end{equation}

\begin{example}
在 $[0,1]$ 上把 $\{1,x\}$ 正交化.
\end{example}
\begin{solution}
取 $\varphi_1=1$,则 $(\varphi_1,\varphi_1)=1$,$(x,1)=\int_0^1 x\,\mathrm{d}x=1/2$,所以
\begin{equation}
  \varphi_2 = x - \frac12 .
\end{equation}
检验 $\int_0^1 (x-1/2)\,\mathrm{d}x=0$.其范数平方是 $\int_0^1 (x-1/2)^2\,\mathrm{d}x=1/12$,对应的归一化函数是 $\sqrt{12}\,(x-1/2)$.
在对称区间 $[-1,1]$ 上, $(x,1)$ 本来就是零, $1$ 和 $x$ 已经正交,不必再减.
\end{solution}

下面我们来讨论正交函数组的完备性. 设$\varphi_1(x), \varphi_2(x), \cdots, \varphi_n(x)$是正交归一组, 我们可以定义
\begin{equation}
  f(x) = \sum_{i=1}^n c_i \varphi_i(x)
\end{equation}
其中系数$c_i$为
\begin{equation}
  c_i = (f, \varphi_i) = \int f(x) \varphi_i(x) dx
\end{equation}
成为\textbf{展开系数}或\textbf{分量}.

由关系式
\begin{equation}
  \int \left( f(x) - \sum_{i=1}^n c_i \varphi_i(x) \right)^2 dx \geq  0
\end{equation}
展开个平方并逐项积分, 我们可以得到
\begin{equation}
 \int f^2(x) dx +  \sum_{i=1}^n c_i^2  -  2 \sum_{i=1}^n   c_i \int f  \varphi_i  dx \geq 0 
\end{equation}
因此有,
\begin{equation}
  \sum_{i=1}^n c_i^2 \leq N f 
\end{equation}
上式可以令$n \to \infty$, 也可以得到
\begin{equation}
  \sum_{i=1}^{\infty} c_i^2 \leq N f
\end{equation}
这称为\textbf{贝塞尔不等式}(Bessel's inequality). 它对每一个正交归一组都成立, 它证明了展开系数的平方和是收敛有界的.

假如有一正交归一组$\varphi_1(x), \varphi_2(x), \cdots, $, 使得对任意分段连续函数$f(x)$,都有$n$使
平均平方误差
$$
\int \left( f(x) - \sum_{i=1}^n c_i \varphi_i(x) \right)^2 dx 
$$
小于任意给定正数$\epsilon$, 那么我们称$\varphi_1(x), \varphi_2(x), \cdots$是\textbf{完备的}. 对于完备的正交归一组而言, 贝塞尔不等式
对任意函数$f(x)$都成为等式:
\begin{equation}
  \sum_{i=1}^{\infty} c_i^2 = N f
  \label{eq:completeness}
\end{equation}
这个关系称为``完备性关系''. 更一般的形式为
\begin{equation}
  \sum_{i=1}^{\infty} c_i d_i = (f, g) \quad \text{其中} \quad  c_i = (f, \varphi_i) , d_i = (g, \varphi_i)
\end{equation}
可以验证, 对$f+g$ 应用完备性关系~\ref{eq:completeness},
\begin{equation}
  N (f + g) = Nf  + Ng + 2 (f, g) = \sum_{i = 1}^{\infty} (c_i + d_i)^2 = \sum_{i=1}^{\infty} (c_i^2 + d_i^2 + 2  c_i d_i) 
\end{equation}
并减去相应的$f,g$的范数.

%  该不等式说明了展开系数的平方和是有界的, 但并不能保证展开系数的平方和等于函数的模平方. 也就是说, 我们并不能保证

\subsection{傅里叶级数的引入}
\label{subsec:fourier_series}
Fourier曾试图解决一个问题,在此过程中发展了傅里叶级数的概念.此问题处理长度为$\ell$的一维均匀棒的热传播过程,即
给定初始温度分布,求解在$t$时刻,$x$处的温度,$T(x,t)$.该温度场满足微分方程
$$
\frac{\partial^2}{\partial x^2} T(x,t) = \frac{\partial}{\partial t} T(x,t), \quad  T(0,t) = T(\ell,t) = 0 .  
$$
该微分方程的一特殊解可以由分离变量法得到,具体方法后面会学习.令$T(x,t) = f(x) g(t)$,代入该方程有
$$
  f''(x) g(t) = f(x) g'(t) ,
$$
两边同除以$f(x)g(t)$得到
$$
\frac{f''(x)}{f(x)} = \frac{g'(t)}{g(t)}   .
$$
由于左式是与$t$无关的,而右边是与$x$无关的,于是唯一可能相等的情况为二者皆为同一个常数$c$.
这时候若$c>0$,可知$\lim_{t\to \infty} g(t) = \infty$,因此我们令$c=-k^2$.
那么可以知道
$$
f(x) =  A \sin{k x} + B \cos{k x},    
$$
由边界条件确定$B = 0$,且$k = \frac{n\pi}{\ell}$.求解$g(t)$得 $g(t) = e^{-k^2 t}$.
该方程的解可以表示为
$$
T(x,t) = A \sin{\left( \frac{n\pi}{\ell} x \right)} e^{-\frac{n^2\pi^2}{\ell^2} t}  
$$
其中$A$由其他条件确定.
可以知道$t=0$时,$T$为一个正弦函数,节点数由$n$决定.但是,更普遍的情况是起始时刻的温度分布并不是正弦函数.
那该如何决定稍后某时刻$t$的分布呢?傅里叶发现我们并不需要对此问题反反复复求解.
对于线性微分方程可以知道若$T_1(x,t)$和$T_2(x,t)$满足微分方程
$$
    \frac{\partial^2}{\partial x^2} T(x,t) = \frac{\partial}{\partial t} T(x,t)
$$
那么二者的任意线性组合
$$
T_3 = \alpha T_1  + \beta T_2  
$$
也是其解.
傅里叶得到结论,对于任意一种光滑的温度分布,其解总可以写成一系列有着不同振幅和模式的正弦函数的线性叠加,即
$$
  T(x,t) = \sum_{n=1}^{\infty} A_n \sin {\left( \frac{n\pi}{\ell} x \right)} e^{-\frac{n^2\pi^2}{\ell^2} t} .     
$$
若初始温度分布函数为$f(x) = T(x,0)$,那么
$$
  f(x) =      \sum_{n=1}^{\infty} A_n \sin {\left( \frac{n\pi}{\ell} x \right)}
$$
上式什么样的函数可以表示为正弦和余弦函数的叠加呢?每个模式的振幅$A_n$该如何计算呢?
这就利用了正余弦函数的特征.

记
\begin{equation}
 (f, g) = \int_0^{2\pi} f^{*}(x) g(x) dx\,,
\end{equation}

可以验证,积分区间取 $[0,2\pi]$ 时
$$
(\sin{ n x }, \sin{m x })  =  (\cos{ n x } \cos {m x })  =  \delta_{mn}\pi  \; (m\neq 0, n\neq 0)
$$

$$
(\sin{ n x },  \cos{m x })  =  0
$$

\begin{example}
用积化和差核对 $[0,2\pi]$ 上的正交性,再换到热棒的区间 $[0,\ell]$.
\end{example}
\begin{solution}
$\sin A\sin B=\frac12\bigl[\cos(A-B)-\cos(A+B)\bigr]$.$m,n$ 是正整数且 $m\neq n$ 时,两个余弦都在 $[0,2\pi]$ 上积满整数个周期,积分是零.$m=n$ 时
\begin{equation}
  \int_0^{2\pi}\frac12\bigl(1-\cos 2nx\bigr)\,\mathrm{d}x=\pi.
\end{equation}
余弦与余弦同样.正弦与余弦化成 $\sin(n+m)x$ 和 $\sin(n-m)x$,在整周期上的积分是零.

热棒的区间是 $[0,\ell]$.令 $\theta=\pi x/\ell$,则 $\mathrm{d}x=(\ell/\pi)\,\mathrm{d}\theta$,$x$ 从 $0$ 到 $\ell$ 时 $\theta$ 从 $0$ 到 $\pi$.
\begin{equation}
  \int_0^{\ell}
    \sin\Bigl(\frac{n\pi x}{\ell}\Bigr)
    \sin\Bigl(\frac{m\pi x}{\ell}\Bigr)
    \,\mathrm{d}x
  =\frac{\ell}{\pi}\int_0^{\pi}\sin(n\theta)\sin(m\theta)\,\mathrm{d}\theta
  =\frac{\ell}{2}\,\delta_{mn}.
\end{equation}
$[0,\pi]$ 上的积分是 $[0,2\pi]$ 上的一半,所以范数从 $\pi$ 变成 $\ell/2$.用 $\sin(m\pi x/\ell)$ 去乘正弦级数再积分,只留下 $A_m\cdot\ell/2$.
\end{solution}

因此振幅系数
\begin{equation}
  A_n = \frac{2}{\ell} \int_0^{\ell} f(x) \sin{ \left( \frac{n\pi}{\ell} x \right) } dx 
\end{equation}
$A_n$称为傅里叶级数的\textbf{系数}.若所有系数都求得后,我们就得到了$f(x)$.原则上,所有的平滑的单值函数并没有什么问题,
但似乎缺乏数学的严谨.自然,越多的级数项对函数的逼近也就越好.

\subsection{傅里叶定理}
\label{subsec:fourier_theorem}
对于任意一个以$2\ell$为周期的函数,
$$
   f(x + 2\ell) = f(x),  
$$
我们可以通过三角函数族进行级数展开
\begin{equation}
  f(x) = \frac{a_0}{2} + \sum_{n=1}^{\infty} \left[ a_n \cos{ \frac{n\pi}{\ell} x } + b_n \sin{ \frac{n\pi}{\ell} x } \right] 
\end{equation}
利用三角函数族的正交性,可以得到展开系数
\begin{align}
  a_n = \frac{1}{\ell} \int_{-\ell}^{\ell} f(x) \cos {  \left( \frac{n\pi}{\ell} x \right) } dx
  \\
  b_n = \frac{1}{\ell} \int_{-\ell}^{\ell} f(x) \sin {  \left( \frac{n\pi}{\ell} x \right) } dx
\end{align}
这里的积分上下限为一个周期$2\ell$,也可以为$[0,2\ell]$.常数项写成 $a_0/2$,是为了让 $n=0$ 仍用同一个积分.$\int_{-\ell}^{\ell}1\,\mathrm{d}x=2\ell$,而 $n\geq 1$ 时 $\int_{-\ell}^{\ell}\cos^2(n\pi x/\ell)\,\mathrm{d}x=\ell$.展开式乘 $1$ 再积分,左边是 $\int_{-\ell}^{\ell}f\,\mathrm{d}x$,右边是 $(a_0/2)\cdot 2\ell$,所以 $a_0$ 正好等于 $n=0$ 的余弦公式.级数里的常数项 $a_0/2$ 就是 $f$ 在一个周期上的平均值.
上述傅里叶级数展开的成立条件被称为\textbf{狄里希利条件} (Dirichlet conditions),即$f(x)$在$\left[-\ell, \ell\right]$区间内只有有限个间断点,且每个周期内有有限个极值点.满足
这两个条件的函数称为\textbf{分段平常}.
满足狄里希利条件的函数$f(x)$在点$x_0$处间断,那么其傅里叶级数在此点的值为该函数左右值的算术平均
\begin{equation}
  \textrm{级数和}(x_0) = \lim_{\epsilon \to 0} \half \left[
     f(x_0 + \epsilon) + f(x_0  - \epsilon) \right]
\end{equation}

对于奇函数和偶函数的傅里叶展开,不难发现,奇函数的展开为

\begin{equation}
  f(x) = \sum_{n=1}^{\infty} b_n \sin {  \left( \frac{n\pi}{\ell} x \right) }
\end{equation}
称为\textbf{傅里叶正弦级数}.
而偶函数的展开为
\begin{equation}
  f(x) = \frac{a_0}{2} + \sum_{n=1}^{\infty}  a_n \cos{  \left( \frac{n\pi}{\ell} x \right) }
\end{equation}
称为\textbf{傅里叶余弦级数}.

只在 $(0,\ell)$ 上给定的函数可以有很多种延拓,它们在 $(0,\ell)$ 上都代表原来的 $f$.边界条件常常把延拓定下来.要求 $f(0)=f(\ell)=0$ 时,把 $f$ 延拓成以 $2\ell$ 为周期的奇函数:在 $(0,\ell)$ 上 $g=f$,而 $-\ell<x<0$ 时 $g(x)=-f(-x)$.由周期性,$g(\ell+\epsilon)=g(-\ell+\epsilon)$;再由奇性,$g(-\ell+\epsilon)=-g(\ell-\epsilon)$.因此 $g$ 关于 $x=\ell$ 也是奇的.连续延拓只能有 $g(0)=g(\ell)=0$.正弦级数在这两个端点自动为零,系数是
\begin{equation}
  b_n = \frac{2}{\ell}\int_0^{\ell} f(x)\sin\Bigl(\frac{n\pi x}{\ell}\Bigr)\,\mathrm{d}x.
\end{equation}
热棒两端保持零度,用的就是这个奇延拓.

要求 $f'(0)=f'(\ell)=0$ 时,改成偶延拓,$g(-x)=g(x)$.偶函数的导数是奇函数,左右导数在 $x=0$ 处相等当且仅当 $f'(0)=0$.把原点移到 $x=\ell$,同样得到 $f'(\ell)=0$.余弦级数的导数是正弦级数,在两端自动为零.系数是
\begin{equation}
  a_n = \frac{2}{\ell}\int_0^{\ell} f(x)\cos\Bigl(\frac{n\pi x}{\ell}\Bigr)\,\mathrm{d}x.
\end{equation}
只在 $(0,\ell)$ 上给定、再选一种延拓,得到的正弦级数或余弦级数称为\idx{半幅级数}{banfujishu}.同一定义可以有两个级数,它们在 $(0,\ell)$ 内都等于 $f$,在区间外代表不同的延拓.

\begin{example}
$f(x)=x$ 定义在 $(0,\ell)$ 上.分别求它的正弦级数和余弦级数.
\end{example}
\begin{solution}
奇延拓后 $b_n=(2/\ell)\int_0^{\ell}x\sin(n\pi x/\ell)\,\mathrm{d}x$.分部积分,$\int_0^{\ell}\cos(n\pi x/\ell)\,\mathrm{d}x=0$,只剩边界项,
\begin{equation}
  b_n=\frac{2\ell}{n\pi}(-1)^{n+1},
  \qquad
  x=\frac{2\ell}{\pi}\sum_{n=1}^{\infty}\frac{(-1)^{n+1}}{n}\sin\Bigl(\frac{n\pi x}{\ell}\Bigr),
  \quad 0<x<\ell.
\end{equation}
端点 $x=0$ 和 $x=\ell$ 上级数和是 $0$,与 $f$ 在端点的值一致.把这个奇延拓写到 $(-\ell,\ell)$ 上,就是后面的锯齿函数.

偶延拓是 $|x|$.$a_0=(2/\ell)\int_0^{\ell}x\,\mathrm{d}x=\ell$,常数项是 $\ell/2$.$n\geq 1$ 时分部积分,
\begin{equation}
  \int_0^{\ell}x\cos\Bigl(\frac{n\pi x}{\ell}\Bigr)\,\mathrm{d}x
  =\frac{\ell^2}{n^2\pi^2}\bigl[(-1)^n-1\bigr],
\end{equation}
边界上的正弦项是零.因此偶数 $n$ 的 $a_n$ 是零,奇数 $n=2k+1$ 给出 $a_n=-4\ell/(n^2\pi^2)$,
\begin{equation}
  x=\frac{\ell}{2}-\frac{4\ell}{\pi^2}\sum_{k=0}^{\infty}
    \frac{1}{(2k+1)^2}\cos\Bigl(\frac{(2k+1)\pi x}{\ell}\Bigr),
  \quad 0<x<\ell.
\end{equation}
令 $x=0$,左边是 $0$,得到 $\sum_{k=0}^{\infty}1/(2k+1)^2=\pi^2/8$.
\end{solution}

\begin{example}
$f(x)=e^{ax}$ ($a\neq 0$) 在 $(-\pi,\pi)$ 上的全级数.周期取 $2\pi$,即 $\ell=\pi$.
\end{example}
\begin{solution}
\begin{equation}
  a_0=\frac{1}{\pi}\int_{-\pi}^{\pi}e^{ax}\,\mathrm{d}x=\frac{2\sinh(a\pi)}{a\pi},
\end{equation}
常数项是 $a_0/2=\sinh(a\pi)/(a\pi)$.对 $n\geq 1$,
\begin{equation}
  \int e^{ax}\cos(nx)\,\mathrm{d}x
  =\frac{e^{ax}\bigl(a\cos(nx)+n\sin(nx)\bigr)}{a^2+n^2}.
\end{equation}
$\sin(n\pi)=0$,$\cos(n\pi)=(-1)^n$,从 $-\pi$ 积到 $\pi$ 得到 $2a(-1)^n\sinh(a\pi)/(a^2+n^2)$.除以 $\pi$,
\begin{equation}
  a_n=\frac{2a(-1)^n\sinh(a\pi)}{\pi(a^2+n^2)}.
\end{equation}
正弦积分
\begin{equation}
  \int e^{ax}\sin(nx)\,\mathrm{d}x
  =\frac{e^{ax}\bigl(a\sin(nx)-n\cos(nx)\bigr)}{a^2+n^2}
\end{equation}
同样只留下余弦边界项,
\begin{equation}
  b_n=\frac{2n(-1)^{n+1}\sinh(a\pi)}{\pi(a^2+n^2)}.
\end{equation}
$a_n$ 和 $b_n$ 都不为零:指数函数既不是奇函数也不是偶函数.系数随 $1/n^2$ 和 $1/n$ 衰减,间断在周期端点 $x=\pm\pi$,那里级数和是两端极限的平均 $\cosh(a\pi)$.
\end{solution}

对于复指数的展开,不难发现可以写成
\begin{equation}
  f(x) = \sum_{n=-\infty}^{\infty} c_n e^{\imath \frac{n\pi}{\ell} x},
\end{equation}
其中
$c_n = \half(a_n - \imath b_n)$, $c_{-n} = \half (a_n + \imath b_n)$, $n>0$, $c_0 = \half a_0$.
写出来为
\begin{equation}
  c_n = \frac{1}{2\ell} \int_{-\ell}^{\ell} f(x) e^{-\imath \frac{n\pi}{\ell} x} dx ,
\end{equation}
尽管$f(x)$是实数,但其傅里叶系数却可能是复数,还可以看出$c_{-n} = c_{n}^*$.
复指数函数族也是正交的
$$
  \braket{ e^{\imath \frac{n\pi}{\ell} x } | e^{\imath \frac{m\pi}{\ell} x } } = 2\ell \delta_{mn} .
$$

\begin{example}
对于锯齿函数
$$  
f(x)= \begin{cases}
  x,   &  0 < x \leq  \ell 
  \\
  x - 2 \ell,   & \ell < x \leq 2\ell
\end{cases}
$$
求其傅里叶级数.
\end{example}
\begin{solution}
把一个周期挪到 $(-\ell,\ell)$,这个函数就是奇函数 $f(x)=x$,间断点在 $\pm\ell$.因此 $a_n=0$,并且
\begin{equation}
  b_n = \frac{2}{\ell}\int_0^{\ell} x\sin\Bigl(\frac{n\pi x}{\ell}\Bigr)\,\mathrm{d}x .
\end{equation}
分部积分,令 $u=x$,$\mathrm{d}v=\sin(n\pi x/\ell)\,\mathrm{d}x$,边界项是 $-\ell^2(-1)^n/(n\pi)$,剩余的余弦积分等于零,所以
\begin{equation}
  b_n = \frac{2\ell}{n\pi}(-1)^{n+1}.
\end{equation}
因此,该级数为
\begin{align}
f(x) &=  \sum_{n=1}^{\infty} \frac{2 \ell }{n\pi} (-1)^{n+1} \sin {  \left( \frac{n\pi}{\ell} x \right) }  \nonumber
\\
 & =\frac{2\ell}{\pi} \left[ \sin{\frac{\pi x}{\ell} } - \frac{1}{2}\sin{\frac{2\pi x}{\ell} } 
  +  \frac{1}{3}\sin{\frac{3\pi x}{\ell} } + \cdots + \frac{(-1)^{n+1}}{n}\sin{\frac{n\pi x}{\ell} }
 \right] \nonumber
.
\end{align}
容易验证
$$
  f(0) = 0; f(\ell) = 0;
$$
当$x=\ell/2$, 
$$
f(\ell/2) = \ell/2 =   \frac{2\ell}{\pi} \left[ 1 - 0 - \frac{1}{3} + \frac{1}{5} - 0 -\frac{1}{7}+ \cdots \right]
$$
因此,我们得到莱布尼兹等式,即以下等式
$$
\frac{\pi}{4} = 1-\frac{1}{3} + \frac{1}{5} - \frac{1}{7} + \cdots = \sum_{n=0} \frac{(-1)^n}{2n + 1} .
$$
上式可以通过等式
$\int_0^1 \frac{d x}{1+x^2}=\left.\tan ^{-1} x\right|_0 ^1=\frac{\pi}{4}$,和被积函数的级数展开验证.
其实,这里的锯齿函数可以直接用$f(x) = x, -\ell < x < \ell$来表示,相应的间断点则被移到$\pm \ell$处.
\end{solution}


\subsection{能量定理}
\idx{帕塞瓦尔等式}{paseiwaer}
由函数的复数傅里叶级数展开可以得到
\begin{equation}
  E = \frac{1}{2\ell} \int_{-\ell}^{\ell} dx f(x)^* f(x) 
 = \frac{1}{2\ell} \int_{-\ell}^{\ell} dx  \sum c_n e^{\imath \frac{n\pi}{\ell} x} \sum c_m e^{-\imath \frac{m\pi}{\ell} x}
 = \sum |c_n|^2
\end{equation}
利用前面的等式$c_n = \half(a_n - \imath b_n)$,  $c_{-n} = \half (a_n + \imath b_n)$, $n>0$, $c_0 = \half a_0$ 可得
\begin{equation}
  E = \frac{a_0^2}{4}+\sum_{n=1}^{\infty}\left(\frac{a_n^2}{2}+\frac{b_n^2}{2}\right) .
\end{equation}
称为\textbf{Parseval等式}(Parseval's identity).
其物理意义指的是,一复杂波形分解为许多波幅为$c_n$的正(余)弦波的叠加,该波所携带的总能量为各个成分波强度的求和.

\begin{example}
矩形波.取 $\ell=1$,
\begin{equation}
  f(x)=\begin{cases}
    -1, & -1<x<0,\\
    +1, & 0<x<1,
  \end{cases}
\end{equation}
再以 $2$ 为周期延拓.
\end{example}
\begin{solution}
$f$ 是奇函数,$a_n=0$.
\begin{equation}
  b_n=2\int_0^1\sin(n\pi x)\,\mathrm{d}x
  =\frac{2\bigl(1-(-1)^n\bigr)}{n\pi}.
\end{equation}
偶数 $n$ 的系数是零,奇数 $n=2k+1$ 的系数是 $4/(n\pi)$,
\begin{equation}
  f(x)=\frac{4}{\pi}\sum_{k=0}^{\infty}
    \frac{\sin\bigl((2k+1)\pi x\bigr)}{2k+1}.
\end{equation}
$x=0$ 是跳跃点,左右极限是 $-1$ 和 $1$,级数和是它们的平均 $0$,每一项在这里确实是零.$x=1/2$ 处 $f=1$,代入得到莱布尼兹级数 $\pi/4=1-1/3+1/5-\cdots$.跳跃附近部分和会冲过阶跃一段固定的比例,再在两侧振荡.项数增加时,这段过冲并不变小,只是挤向间断点.平均平方误差仍然趋于零:逐点图像和平均平方逼近回答的是两件不同的事.

$|f|^2=1$,所以 $E=1/(2\ell)\int_{-\ell}^{\ell}|f|^2\,\mathrm{d}x=1$.Parseval 等式只留下奇数项,
\begin{equation}
  1=\sum_{k=0}^{\infty}\frac12\Bigl(\frac{4}{(2k+1)\pi}\Bigr)^2
  =\frac{8}{\pi^2}\sum_{k=0}^{\infty}\frac{1}{(2k+1)^2}.
\end{equation}
因此 $\sum_{k=0}^{\infty}1/(2k+1)^2=\pi^2/8$.这和半幅余弦级数在 $x=0$ 处得到的是同一个和.
\end{solution}

\begin{example}
$f(x)=x^2$ 在 $[-\pi,\pi]$ 上,周期 $2\pi$.
\end{example}
\begin{solution}
这是偶函数,$b_n=0$.
\begin{equation}
  a_0=\frac{1}{\pi}\int_{-\pi}^{\pi}x^2\,\mathrm{d}x=\frac{2\pi^2}{3},
\end{equation}
常数项 $a_0/2=\pi^2/3$.分部积分两次,边界上的 $\sin(nx)$ 是零,
\begin{equation}
  \int_0^{\pi}x^2\cos(nx)\,\mathrm{d}x=\frac{2\pi(-1)^n}{n^2},
  \qquad
  a_n=\frac{2}{\pi}\cdot\frac{2\pi(-1)^n}{n^2}=\frac{4(-1)^n}{n^2}.
\end{equation}
\begin{equation}
  x^2=\frac{\pi^2}{3}+4\sum_{n=1}^{\infty}\frac{(-1)^n}{n^2}\cos(nx),
  \qquad -\pi\leq x\leq\pi.
\end{equation}
$x=\pi$ 处函数连续,$x^2=\pi^2$,$\cos(n\pi)=(-1)^n$,
\begin{equation}
  \pi^2=\frac{\pi^2}{3}+4\sum_{n=1}^{\infty}\frac{1}{n^2},
\end{equation}
所以 $\sum_{n=1}^{\infty}1/n^2=\pi^2/6$.

再用 Parseval 等式核对.左边
\begin{equation}
  E=\frac{1}{2\pi}\int_{-\pi}^{\pi}x^4\,\mathrm{d}x=\frac{\pi^4}{5}.
\end{equation}
右边是 $a_0^2/4+\sum a_n^2/2=(2\pi^2/3)^2/4+\sum_{n=1}^{\infty}\frac12\bigl(4/n^2\bigr)^2=\pi^4/9+8\sum 1/n^4$.这个等式给出的是 $\sum 1/n^4$,不是 $\sum 1/n^2$.$\pi^2/6$ 来自把级数代回 $x=\pi$,两边各自算完再比较.
\end{solution}

\subsection{三角系的完备性}
前面把完备说成:平均平方误差可以小于任意正数.这里对分段平常的周期函数把这件事证完.先作换元 $\theta=\pi x/\ell$,把区间 $[-\ell,\ell]$ 化到 $[-\pi,\pi]$.两边的积分只差常数 $\ell/\pi$,所以只要在 $[-\pi,\pi]$ 上证明即可.此时
\begin{equation}
  a_n = \frac{1}{\pi}\int_{-\pi}^{\pi} f(x)\cos(nx)\,\mathrm{d}x,
  \qquad
  b_n = \frac{1}{\pi}\int_{-\pi}^{\pi} f(x)\sin(nx)\,\mathrm{d}x.
\end{equation}

设 $s_n$ 是 $f$ 的 $n$ 阶傅里叶部分和,$T$ 是次数不超过 $n$ 的三角多项式,$r=T-s_n$.由系数定义,$\int_{-\pi}^{\pi}(f-s_n)\cos(kx)\,\mathrm{d}x=0$,正弦同样,所以 $\int(f-s_n)r\,\mathrm{d}x=0$.于是
\begin{equation}
  \int_{-\pi}^{\pi}(f-T)^2\,\mathrm{d}x
  = \int_{-\pi}^{\pi}(f-s_n)^2\,\mathrm{d}x + \int_{-\pi}^{\pi} r^2\,\mathrm{d}x.
\end{equation}
第二项非负,因此在这个次数之内,$s_n$ 的平均平方误差最小.把平方展开,并用 $\int_{-\pi}^{\pi}\cos^2(kx)\,\mathrm{d}x=\pi$($k\geq 1$)和 $\int_{-\pi}^{\pi}1\,\mathrm{d}x=2\pi$,得到
\begin{equation}
  \int_{-\pi}^{\pi}(f-s_n)^2\,\mathrm{d}x
  = \int_{-\pi}^{\pi} f^2\,\mathrm{d}x
    - \pi\Biggl(\frac{a_0^2}{2}+\sum_{k=1}^{n}(a_k^2+b_k^2)\Biggr).
\end{equation}
左边非负,这就是三角系上的贝塞尔不等式.要证明完备,还要造出一组三角多项式,把这个差压到任意小.

部分和可以写成卷积.记\textbf{狄里希利核}
\begin{equation}
  D_n(\theta) = \half + \sum_{k=1}^{n}\cos(k\theta)
  = \frac{\sin\bigl((n+\half)\theta\bigr)}{2\sin(\theta/2)}.
\end{equation}
第二个等号用 $2\sin(\theta/2)$ 去乘求和式,积化和差后中间项相消.于是
\begin{equation}
  s_n(x) = \frac{1}{\pi}\int_{-\pi}^{\pi} f(t)\,D_n(t-x)\,\mathrm{d}t,
\end{equation}
并且 $\int_{-\pi}^{\pi} D_n(\theta)\,\mathrm{d}\theta=\pi$.

把前 $n+1$ 个狄里希利核取平均,得到\textbf{费耶核}
\begin{equation}
  F_n(\theta)
  = \frac{1}{n+1}\sum_{m=0}^{n} D_m(\theta)
  = \frac{1}{2(n+1)}
    \Biggl(\frac{\sin((n+1)\theta/2)}{\sin(\theta/2)}\Biggr)^2.
\end{equation}
它非负,积分仍是 $\pi$.对应的\textbf{费耶平均} $\sigma_n=(s_0+\cdots+s_n)/(n+1)$ 仍是次数不超过 $n$ 的三角多项式,并且
\begin{equation}
  \sigma_n(x)-f(x)
  = \frac{1}{\pi}\int_{-\pi}^{\pi}\bigl[f(x-t)-f(x)\bigr] F_n(t)\,\mathrm{d}t.
\end{equation}

先设 $f$ 连续,并且 $f(-\pi)=f(\pi)$,从而在圆周上一致连续.给定 $\epsilon>0$,取 $\delta$ 使 $|t|<\delta$ 时 $|f(x-t)-f(x)|<\epsilon$.积分拆成 $|t|<\delta$ 和 $|t|\geq\delta$.$F_n\geq 0$,峰内的贡献小于 $\epsilon$.峰外 $|\sin(t/2)|\geq\sin(\delta/2)$,所以 $F_n$ 不超过 $1/\bigl(2(n+1)\sin^2(\delta/2)\bigr)$,$n$ 大时峰外的贡献可以小于 $\epsilon$.因此 $\sigma_n$ 一致收敛到 $f$,
\begin{equation}
  \int_{-\pi}^{\pi}(f-\sigma_n)^2\,\mathrm{d}x
\end{equation}
可以任意小.$s_n$ 的平均平方误差不超过 $\sigma_n$ 的,所以连续且周期闭合的函数已经可以用三角多项式在均方意义下逼近.

若 $f$ 只是分段平常,一个周期里只有有限个跳跃.在每个跳跃处用一段宽度为 $\eta$ 的直线把左右值接上,端点 $f(-\pi)$ 与 $f(\pi)$ 同样接上,得到连续的周期函数 $h$.这些小区间上 $|f-h|$ 有界,$\eta$ 足够小时 $\int(f-h)^2$ 可以任意小.再对 $h$ 取三角多项式 $T$,使 $\int(h-T)^2$ 任意小.由 $(A+B)^2\leq 2A^2+2B^2$,
\begin{equation}
  \int(f-T)^2\,\mathrm{d}x
  \leq 2\int(f-h)^2\,\mathrm{d}x + 2\int(h-T)^2\,\mathrm{d}x
\end{equation}
也可以任意小.$n$ 不小于 $T$ 的次数时,$\int(f-s_n)^2\leq\int(f-T)^2$.平均平方误差可以任意小,三角函数族完备.回到 $[-\ell,\ell]$,贝塞尔不等式里的余项趋于零,就得到正文中的 Parseval 等式.狄里希利条件另外保证间断点处级数和等于左右值的平均;这里证明的是平均平方意义下的完备.
\section{傅里叶变换}
\label{sec:fourier_transform}

\subsection{傅里叶变换}
\label{subsec:fourier_transform}
前面我们讨论了周期性体系的傅里叶级数展开.现在我们希望研究非周期性函数的傅里叶展开问题.
假设$f(x)$是定义在$-\infty < x < \infty$的实函数,我们依旧对该函数进行傅里叶级数展开,
该级数必须理解为周期$2\ell$扩展至$\infty$的极限情况.
由于三角函数族的变量为
\[
 \frac{n\pi x}{\ell}    
\]
引入非连续变量
\[\omega_n = \frac{n\pi} {\ell}, \quad k = 0,1,2,\cdots
    \]
可知$\omega_n = n \omega_0$, 其中$\omega_0 = \frac{\pi}{\ell}$为基础频率(Fundamental Frequency),
\[\Delta \omega_n = \omega_n - \omega_{n-1} = \omega_0 . \] 
有
\begin{equation}
    f(x) = \lim_{\ell\to \infty} \sum_{n=-\infty}^{\infty} \left[
        c_n e^{\imath \omega_n x}
        % \frac{a_0}{2} + \sum_{n=1}^{\infty} \left( a_n \cos \omega_n x + b_n \sin \omega_n x \right)
    \right]
\end{equation}
其中系数
\begin{equation}
    c_n = \frac{1}{2\ell} \int_{-\ell}^{\ell} f(x) e^{-\imath \omega_n x} dx 
\end{equation}
取$\ell\to \infty$的极限,则
\[
  \sum_n \to  \frac{ \ell}{\pi}  \int d\omega 
\]
傅里叶系数记为$F(\omega)$,
有
\begin{equation}
    f(x) = \int_{-\infty}^{\infty} d \omega F(\omega) e^{\imath \omega x}, % \quad F (\omega) \equiv 2\ell c_n
\end{equation}
上式的积分称为\textbf{傅里叶积分}(Fourier integral),是$f(x)$的\textbf{傅里叶变换}(Fourier transform).
其中
\begin{equation}
    F (\omega) \equiv \lim_{\ell\to \infty} \frac{ 2\ell c_n}{2\pi}  = \frac{1}{2\pi} \int_{-\infty}^{\infty} e^{-\imath \omega x} f(x) dx,
\end{equation}
这被称为\textbf{逆傅里叶变换}(inverse Fourier transform). 这里积分前的系数并不要紧,更关键的是积分关于$\omega$的形式.

傅里叶积分存在的条件由\textbf{傅里叶积分定理}判定: 
若函数 $f(x)$ 在区间 $(-\infty, \infty)$ 上满足条件: $(1) f(x)$ 在 任一有限区间上满足狄里希利条件; (2) $f(x)$ 在 $(-\infty, \infty)$ 
上绝对可积 (即 $\int_{-\infty}^{\infty}|f(x)| dx$ 收敛), 则 $f(x)$ 可表示成傅里叶积分, 且
$$
\textrm{傅里叶积分值} = \lim_{\epsilon\to 0} \half \left[ f(x+\epsilon) + f(x-\epsilon) . \right]
$$
% $F (\omega) \equiv 2\ell c_n$.

% \begin{align}
%     a_n =  \lim_{\ell\to \infty} \frac{1}{\ell} \int_{-\ell}^{\ell} f(x) \cos {  \omega_n x  } dx
%     \\
%     b_n = \frac{1}{\ell} \int_{-\ell}^{\ell} f(x) \sin {   \omega_n x } dx
% \end{align}

正余弦形式利用 $c_n$ 与 $a_n$、$b_n$ 的关系得到.
上一节已有 ($n>0$)
\[
  c_n=\half(a_n-\imath b_n),
  \qquad
  c_{-n}=\half(a_n+\imath b_n),
\]
亦即 $a_n=c_n+c_{-n}$, $b_n=\imath(c_n-c_{-n})$.
把 $\pm n$ 两项合并:
\[
  c_n e^{\imath\omega_n x}+c_{-n}e^{-\imath\omega_n x}
  =a_n\cos(\omega_n x)+b_n\sin(\omega_n x).
\]
又 $F(\omega)=\lim(2\ell c_n)/(2\pi)=\lim c_n/\Delta\omega$, 相应定义
\[
  A(\omega)=\lim\frac{a_n}{\Delta\omega}=F(\omega)+F(-\omega),
  \quad
  B(\omega)=\lim\frac{b_n}{\Delta\omega}=\imath\bigl[F(\omega)-F(-\omega)\bigr],
\]
或等价地 $F(\omega)=\half\bigl[A(\omega)-\imath B(\omega)\bigr]$.
由 $a_n=(1/\ell)\int f\cos$ 与 $1/\ell=\Delta\omega/\pi$, 取 $\ell\to\infty$ 得
\begin{equation}
    \left\{\begin{array}{l}
    A(\omega)=\frac{1}{\pi} \int_{-\infty}^{\infty} f(t) \cos \omega t\, d t, \\
    B(\omega)=\frac{1}{\pi} \int_{-\infty}^{\infty} f(t) \sin \omega t\, d t .
    \end{array}\right.
\end{equation}
级数中 $\pm n$ 合并后再取极限, 得到正余弦形式
\begin{equation}
    f(x) =\int_{0}^{\infty} A(\omega) \cos {\omega x} d\omega + \int_{0}^{\infty} B(\omega) \sin {\omega x} d\omega.
\end{equation}

上式还可以写成

\begin{equation}
    f(x)=\int_0^{\infty} C(\omega) \cos [\omega x-\varphi(\omega)] d \omega,
\end{equation}

其中
$$
\begin{gathered}
C(\omega)=\sqrt{[A(\omega)]^2+[B(\omega)]^2} \\
\varphi(\omega)=\arctan [B(\omega) / A(\omega)] .
\end{gathered}
$$
$C(\omega)$ 称为 $f(x)$ 的振幅谱, $\varphi(\omega)$ 称为 $f(x)$ 的相位谱.如果$f(x)$是奇函数或偶函数,
对应的级数为\textbf{傅里叶余弦积分}或\textbf{傅里叶正弦积分}.

复数形式的傅里叶积分可以写成对称的形式,
\begin{equation}
    \begin{aligned}
    f(x) & =\frac{1}{\sqrt{2 \pi}} \int_{-\infty}^{\infty} F(\omega)e^{\imath \omega x} d \omega, \\
    F(\omega) & =\frac{1}{\sqrt{2 \pi}} \int_{-\infty}^{\infty} f(x)\left[e^{\imath \omega x}\right]^* d x .
    \end{aligned}
\end{equation}
并常用符号简写为
$$
F(\omega)=\mathcal{F}[f(x)], \quad f(x)=\mathcal{F}^{-1}[F(\omega)] .
$$
$f(x)$ 和 $F(\omega)$ 分别称为傅里叶变换的\textbf{原函数}和\textbf{像函数}.

\begin{example}
求以下函数的复数傅里叶变换.
\begin{enumerate}
    \item $f(x) = e^{-\alpha |x|}, \alpha > 0$;
    \item $f(x) = \delta (x)$;
    \item $f(x) = e^{-\alpha x^2}, \alpha > 0$.
\end{enumerate}
\end{example}
\begin{solution}
\begin{enumerate}
    \item 对$x$分段进行处理,利用对称形式的傅里叶变换有
    \begin{align}
        F(\omega) &=
        \sqrt{\frac{1}{2 \pi}} \int_{-\infty}^0 e^{\alpha x+i \omega t} dx \nonumber
        % \\ &
        +
        \sqrt{\frac{1}{2 \pi}} \int_0^{\infty} e^{-\alpha x+i \omega t} dx  \nonumber
        \\
        &= \sqrt{\frac{1}{2 \pi}} \left[\frac{1}{\alpha + \imath \omega } + \frac{1}{\alpha - \imath \omega } \right] \nonumber
        \\
        &= \sqrt{\frac{1}{2 \pi}} \frac{2\alpha}{\alpha^2 + \omega^2}, \nonumber
    \end{align}
    越大的$\alpha$则$f(x)$越窄,集中在$x=0$附近,也就是越局域,而其傅里叶变换则越弥散.当$f(x)$是偶函数时,傅里叶变换的像函数为实函数.

    \item 代入
    \begin{align}
        F(\omega) &=
        \sqrt{\frac{1}{2 \pi}} \int_{-\infty}^{\infty} \delta(x) e^{+\imath \omega x} d x = \sqrt{\frac{1}{2 \pi}} ,
    \end{align}
    $f(x)$无限局域对应的傅里叶变换为无限弥散.

    \item 类似的,  
    \begin{align}
        F(\omega) &=\frac{1}{\sqrt{2 \pi}} \int_{-\infty}^{\infty} e^{-\alpha x^2} e^{\imath \omega x} d x \nonumber
        % \\ &
        % \\
        % & 
        = \frac{ e^{-\omega^2 / 4\alpha }}{\sqrt{2 \pi}} \int_{-\infty}^{\infty} e^{-\alpha  \left( x - \frac{\imath \omega}{2\alpha} \right)^2}  d x
        \nonumber
        % \\ & 
        \\
        & = \frac{ e^{-\omega^2 / 4\alpha }}{\sqrt{2 \pi}} \int_{-\infty - \imath \omega / 2\alpha}^{\infty - \imath \omega / 2\alpha} e^{-\alpha t^2} dt
        \nonumber
        % \\ &
        \\
        & = \frac{ e^{-\omega^2 / 4\alpha }}{\sqrt{2 \pi}} \sqrt{\frac{\pi}{\alpha}} 
        \nonumber
        % \\ &
        \\
            & = \frac{1}{\sqrt{2\alpha}} e^{- \frac{\omega^2}{4\alpha}} \nonumber .
        % \\ &
    \end{align}
    上式用了变量代换$t = x - \imath \omega /2 \alpha$,高斯函数的像函数还是高斯函数.
\end{enumerate}
\end{solution}


\subsection{傅里叶变换基本性质}
\idx{卷积}{juanji}
傅里叶变换满足以下基本性质.假定$f(x)$的傅里叶变换存在,记为$\mathcal{F}[f(x)] = F(\omega)$.

\begin{enumerate}
    \item 导数定理
    \begin{equation}
        \mathcal{F} [f'(x)] = \imath \omega F(\omega)
    \end{equation}

    \item 积分定理
    \begin{equation}
        \mathcal{F} [ \int^{x} f(x) dx ] = \frac{1}{\imath \omega} F(\omega)
    \end{equation}

    \item 相似性定理
    \begin{equation}
        \mathcal{F} [ f(ax) ] = \frac{1}{|a|} F\Bigl(\frac{\omega}{a}\Bigr),
        \qquad a\neq 0.
    \end{equation}
    绝对值来自雅可比. $a<0$ 时坐标反向,写成 $1/a$ 会丢一个负号.
    \item 延迟定理
    \begin{equation}
        \mathcal{F} [ f(x - x_0 ) ] = e^{-\imath \omega x_0} F(\omega),
    \end{equation}
    \item 位移定理
    \begin{equation}
        \mathcal{F} [ e^{\imath \omega_0 x} f(x) ] = F(\omega - \omega_0),
    \end{equation}
    \item 卷积定理. 卷积定义为
    \begin{equation}
        (f_1 * f_2)(x) = \int_{-\infty}^{\infty} f_1(\alpha)\, f_2(x-\alpha)\, d\alpha .
    \end{equation}
    注意第二个因子的自变量是 $x-\alpha$,不是 $\alpha-x$.在本节的对称约定下,
    \begin{equation}
        \mathcal{F}[f_1 * f_2] = \sqrt{2\pi}\, F_1(\omega)\, F_2(\omega).
    \end{equation}
    因子 $\sqrt{2\pi}$ 来自归一化.若把 $2\pi$ 全部放在逆变换上,这个常数会换成别的数,乘积结构不变.
\end{enumerate}

导数定理可以由分部积分直接看出.边界项在无穷远处消失时,
\begin{equation}
    \int_{-\infty}^{\infty} f'(x)\, e^{-\imath \omega x}\,\mathrm{d}x
    = \imath\omega \int_{-\infty}^{\infty} f(x)\, e^{-\imath \omega x}\,\mathrm{d}x,
\end{equation}
两边再除以 $\sqrt{2\pi}$ 就是 $\mathcal{F}[f']=\imath\omega F$.延迟定理是换元 $u=x-x_0$:指数多出一个 $e^{-\imath\omega x_0}$,积分变回 $F(\omega)$.

\begin{example}
求宽度为 $2a$ 的一维狭缝 $f(x)=1$（$|x|\le a$）的傅里叶变换,并由此说明矩孔.
\end{example}
\begin{solution}
\begin{equation}
    F(\omega)
    = \frac{1}{\sqrt{2\pi}}\int_{-a}^{a} e^{-\imath\omega x}\,\mathrm{d}x
    = \sqrt{\frac{2}{\pi}}\, a\, \frac{\sin(a\omega)}{a\omega}.
\end{equation}
$\omega=0$ 时商式的极限是 $\sqrt{2/\pi}\,a$.缝越宽,中央峰越窄.二维矩孔是两个这样的因子相乘.
\end{solution}

\subsection{高维傅里叶变换}

二维连续函数 $f(x, y)$ 的傅里叶变换定义如下:
设 $f(x, y)$ 是两个独立变量 $x, y$ 的函数, 且在 $\pm \infty$ 上绝对可积, 则定义积分
$$
F\left(k_1, k_2\right)=\frac{1}{2 \pi} \int_{-\infty}^{\infty} \int_{-\infty}^{\infty} f(x, y) e^{-\imath\left(k_1 x+k_2 y\right)} d x d y
$$
为二维连续函数 $f(x, y)$ 的傅里叶变换,并定义
$$
f(x, y)=\int_{-\infty}^{\infty} \int_{-\infty}^{\infty} F\left(k_1, k_2\right) e^{\imath\left(k_1 x+k_2 y\right)} d k_1 d k_2
$$
为 $F\left(k_1, k_2\right)$ 的逆变换.
这是把 $2\pi$ 全部放在正变换上的写法.与一维例子相同的对称约定是两侧各有一个 $1/(2\pi)$:
\begin{equation}
    F(k_1,k_2)=\frac{1}{2\pi}\iint f(x,y)\,e^{-\imath(k_1 x+k_2 y)}\,\mathrm{d}x\,\mathrm{d}y,
\end{equation}
\begin{equation}
    f(x,y)=\frac{1}{2\pi}\iint F(k_1,k_2)\,e^{\imath(k_1 x+k_2 y)}\,\mathrm{d}k_1\,\mathrm{d}k_2.
\end{equation}
两种约定不能在同一次计算里混用.下面的矩孔按对称约定计算.

\begin{example}
    求函数 $f(x, y)= \begin{cases}A, & |x| \leq X,|y| \leq Y \\ 0, & |x|>X,|y|>Y\end{cases}$ 的傅里叶变换 (矩孔费琅和夫衍射).
\end{example}
\begin{solution}
    由傅里叶变换关系
$$
F\left(k_1, k_2\right)=\frac{1}{2 \pi} \int_{-\infty}^{\infty} \int_{-\infty}^{\infty} f(x, y) e^{-i\left(k_1 x+k_2 y\right)} d x d y
$$
有
$$
\begin{aligned}
F\left(k_1, k_2\right) & =\frac{A}{2 \pi} \int_{-X}^X e^{-i k_1 x} d x \int_{-Y}^Y e^{-i k_2 y} d y \\
& =\left.\left.\frac{A}{2 \pi} \frac{1}{-i k_1} e^{-i k_1 x}\right|_{-X} ^X \frac{1}{-i k_2} e^{-i k_2 y}\right|_{-Y} ^Y \\
& =\frac{2A}{\pi} \frac{1}{i 2 k_1}\left(e^{i k_1 X}-e^{-i k_1 X}\right) \frac{1}{i 2 k_2}\left(e^{i k_2 Y}-e^{-i k_2 Y}\right) \\
& =\frac{2A X Y}{\pi} \frac{\sin \left(k_1 X\right)}{k_1 X} \frac{\sin \left(k_2 Y\right)}{k_2 Y}
\end{aligned}
$$
\end{solution}

物理上,$f(x,y)$ 是孔径面上的复振幅透过率:矩形开口内均匀照明取常值 $A$,挡光处为零.
费琅和夫(远场)衍射振幅正是孔径函数的傅里叶变换 $F(k_1,k_2)$,观察屏上的光强正比于 $|F|^2$;
$k_1,k_2$ 对应衍射角方向.
轴上 $F(0,0)=2AXY/\pi$,与孔径面积成正比.
$k_1 X=n\pi$ 或 $k_2 Y=m\pi$($n,m=\pm 1,\pm 2,\ldots$) 处振幅为零,屏上出现互相垂直的暗纹.
孔越窄,中央峰在对应方向越宽,与一维狭缝同一条宽窄互换.
$F$ 又是两个一维狭缝因子的乘积,故矩孔图样是互相垂直的两条单缝衍射结构.

三维时对称约定是两侧各有 $(2\pi)^{-3/2}$:
\begin{equation}
    F(\vec{k})=\frac{1}{(2 \pi)^{3/2}} \iiint_{-\infty}^{\infty} f(\vec{r}) e^{-\imath  \vec{k} \cdot \vec{r}} d^3 \vec{r},
\end{equation}
\begin{equation}
     f(\vec{r})=\frac{1}{(2 \pi)^{3/2}}\iiint_{-\infty}^{\infty} F(\vec{k}) e^{\imath \vec{k} \cdot \vec{r}} d^3 \vec{k}.
\end{equation}
把 $(2\pi)^{-3}$ 只放在一侧是另一种等价约定,不能和上式混用.指数可分离时,积分拆成三个一维积分.

若 $f(\vec{r})=f(r)$ 球对称,则 $F$ 只依赖 $k=|\vec{k}|$.
取极轴沿 $\vec{k}$,角积分给出 $4\pi\sin(kr)/(kr)$,对称约定下
\begin{equation}
    F(k)
    =\sqrt{\frac{2}{\pi}}\,\frac{1}{k}
      \int_0^{\infty} r\,f(r)\sin(kr)\,\mathrm{d}r.
\end{equation}
三维球对称傅里叶变换化成一维正弦变换.

\begin{example}
求汤川势 $f(r)=e^{-\mu r}/r$($\mu>0$) 的三维傅里叶变换,并说明它在散射中的用法.
\end{example}
\begin{solution}
用 $\int_0^{\infty}e^{-\mu r}\sin(kr)\,\mathrm{d}r=k/(k^2+\mu^2)$,得
\begin{equation}
    F(k)=\sqrt{\frac{2}{\pi}}\,\frac{1}{k^2+\mu^2}.
\end{equation}
散射计算里动量空间矩阵元常不带 $(2\pi)^{-3/2}$,直接写
\begin{equation}
    \int\frac{e^{-\mu r}}{r}\,e^{-\imath\vec{q}\cdot\vec{r}}\,\mathrm{d}^3\vec{r}
    =\frac{4\pi}{q^2+\mu^2},
    \qquad \vec{q}=\vec{k}'-\vec{k}.
\end{equation}
这是玻恩近似下一阶散射振幅的势因子.$\mu\to 0$ 给出库仑势的 $4\pi/q^2$,对应卢瑟福公式的动量依赖.
\end{solution}
\section{$\delta$函数}
\idx{狄拉克函数}{dilake}
\label{sub:delta_function}

\subsection{$\delta$函数的概念}

物理学常常要研究一个物理量在空间或时间中分布的密度, 例如质量密度 (通常简称为密度)、电荷密度、每单位时间传递的动量 (即力) 等等. 但是物 理学又常常运用质点、点电荷、瞬时力等抽象模型, 它们不是连续分布于空间 或时间中, 而是集中在空间的某一点或时间的某一瞬时. 它们的密度又该如何 描述呢?
我们知道, 若质量 $m$ 均匀分布在长为 $l$ 的线段 $[-l / 2, l / 2]$ 上, 则其线密 度 $\rho_l(x)$ 可表为
$$
\rho_l(x)=\left\{\begin{array}{ll}
0 & (|x|>l / 2), \\
m / l & (|x| \leqslant l / 2),
\end{array} \text { 即 } \rho_l(x)=\frac{m}{l} \operatorname{rect}\left(\frac{x}{l}\right) .\right.
$$
将 $\rho_l(x)$ 对 $x$ 积分, 则得到总质量
$$
\int_{-\infty}^{\infty} \rho_l(x) d x=\int_{-\frac{l}{2}}^{\frac{l}{2}} \frac{m}{l} d x=m
$$

如果让上述线段的长度 $l \rightarrow 0$, 我们将得到位于坐标原点质量为 $m$ 的质点, 而线密度函数就成为质点的线密度函数. 将它记为 $\rho(x)$, 则
$$
\lim _{l \rightarrow 0} \int_{-\infty}^{\infty} \rho_l(x) d x=\int_{-\infty}^{\infty} \rho(x) d x=m .
$$
若不求积分而先取极限, 则有
$$
\rho(x)=\lim _{l \rightarrow 0} \rho_l(x)=\lim _{l \rightarrow 0} \frac{m}{l} \operatorname{rect}\left(\frac{x}{l}\right)= \begin{cases}0 & (x \neq 0), \\ \infty & (x=0) .\end{cases}
$$

由此可以看出质点线密度分布函数的直观图像, 它在 $x=0$ 处为 $\infty$, 在 $x \neq$ 0 处为零. 它的积分为 $m$.
对于质点、点电荷、瞬时力这类集中于空间某一点或时间的某一瞬时的抽 象模型, 在物理学中引人 $\delta$ 函数以描述其密度:
$$
\begin{gathered}
\delta(x)= \begin{cases}0 & (x \neq 0), \\
\infty & (x=0) ;\end{cases} \\
\int_a^b \delta(x) d x= \begin{cases}0 & (\text{其他情况}), \\
1 & (a<0<b) .\end{cases}
\end{gathered}
$$
\subsection{$\delta$函数的性质}
$\delta$函数的很多性质可以用在实际的计算中,方便实用. 首先它是满足积分为单位1.
$$
\int_{-\infty}^{\infty} \delta(x) d x=1
$$

并且作为被积函数的乘积因子,它具有挑选性特点

$$
\int_{-\infty}^{\infty} f(x) \delta(x-a) d x=f(a).
$$

$\delta(x)$ 是偶函数, 它的导数是奇函数,
$$
\begin{gathered}
\delta(-x)=\delta(x), \\
\delta^{\prime}(-x)=-\delta^{\prime}(x) .
\end{gathered}
$$
对光滑的 $f$,分部积分把导数挪到 $f$ 上,
\begin{equation}
  \int_{-\infty}^{\infty} f(x)\,\delta'(x-a)\,\mathrm{d}x = -f'(a).
\end{equation}
所以 $\delta'$ 挑选的是导数值,并且带一个负号.

对于变量为函数的情况,有
\begin{equation}
    \delta[\varphi(x)]=\sum_k \frac{\delta\left(x-x_k\right)}{\left|\varphi^{\prime}\left(x_k\right)\right|}
\end{equation}
其中$x_k(k=1,2,3, \cdots)$是
$\varphi(x)=0$ 的实根.
特例包括
\begin{equation}
    \begin{gathered}
    \delta(a x)=\frac{\delta(x)}{|a|}, \\
    \delta\left(x^2-a^2\right)=\frac{\delta(x-a)+\delta(x+a)}{2|a|} .
    \end{gathered}
\end{equation}
计算时分母用 $2|a|$.只在零点上 $2|x|$ 才和 $2|a|$ 相等,不要把 $|x|$ 写进恒等式.

\begin{example}
计算 $\displaystyle\int_{-\infty}^{\infty}(x^2+1)\,\delta(x^2-4)\,\mathrm{d}x$.
\end{example}
\begin{solution}
零点是 $x=\pm 2$, $|\varphi'(x)|=|2x|=4$.每个零点贡献 $(x_k^2+1)/4=5/4$,所以积分等于 $5/2$.
\end{solution}

\begin{example}
电荷 $-e$、$2e$、$-e$ 分别位于 $x=-4$、$0$、$\pi$.
\end{example}
\begin{solution}
线密度为 $\rho(x)=-e\,\delta(x+4)+2e\,\delta(x)-e\,\delta(x-\pi)$.
总电荷是系数之和,等于 $0$.区间 $(-3,3)$ 只含 $x=0$,这段上的电荷是 $2e$.
\end{solution}

$\delta$函数是一种广义函数,常常可以看作是以下函数的极限.
\begin{equation}
    \begin{aligned}
    & \delta(x)=\lim _{l \rightarrow 0} \frac{1}{l} \operatorname{rect}\left(\frac{x}{l}\right), \\
    & \delta(x)=\lim _{K \rightarrow \infty} \frac{1}{\pi} \frac{\sin K x}{x}, \\
    & \delta(x)=\lim _{\varepsilon \rightarrow 0} \frac{1}{\pi} \frac{\varepsilon}{\varepsilon^2+x^2} .
    \end{aligned}
\end{equation}

\subsection{其他}
$\delta$ 函数的傅里叶变换: 用复数形式的傅里叶积分表示$\delta$函数,
\begin{equation}
\delta(x)=\int_{-\infty}^{\infty} C(\omega) e^{\imath\omega x} d \omega,
\end{equation}
其中变换
$$
C(\omega)=\frac{1}{2 \pi} \int_{-\infty}^{\infty} \delta(x) e^{-\imath\omega x} d x=\frac{1}{2 \pi} e^{-\imath\omega \cdot 0}=\frac{1}{2 \pi} .
$$
因此,有$\delta$函数的傅里叶积分为
\begin{equation}
\delta(x)=\frac{1}{2 \pi} \int_{-\infty}^{\infty} e^{i \omega x} d \omega .
\end{equation}

多维的 $\delta$ 函数:
$$
\begin{aligned}
& \br=(x,y,z), \quad \delta (\br)= \begin{cases}0,  & \br \neq 0) \\
\infty, & (\br=0)\end{cases} \\
& \iiint_V \delta(\br)d^3 \br=\iiint_{-\infty}^{\infty} \delta(\br) dx dy dz=1
\end{aligned}
$$
在直角坐标系中, 这样的三维 $\delta$ 函数往往用三个一维 $\delta$ 函数的乘积表示:
\begin{equation}
    \delta(\br)=\delta(x) \delta(y) \delta(z),
\end{equation}
场点不在原点、也不在球坐标的极轴上时,分母是体积元的雅可比:
\begin{align}
    \delta(\br-\brp)
    &=\frac{\delta(\rho-\rho')\,\delta(\varphi-\varphi')\,\delta(z-z')}{\rho},\\
    \delta(\br-\brp)
    &=\frac{\delta(r-r')\,\delta(\theta-\theta')\,\delta(\varphi-\varphi')}{r^2\sin\theta}.
\end{align}
上一行是柱坐标,下一行是球坐标.这样对 $\delta$ 积分才和体积元配成 $1$.原点处的柱坐标表达式不能省掉 $\delta(\rho)$.
\section{拉普拉斯变换与反变换}
\idx{拉普拉斯变换}{lapulasibianhuan}
\label{sec:laplace_transform_and_inverse_transform}
\subsection{拉普拉斯变换}
\label{subsec:laplace_transform}

傅里叶变换在分析信号的频谱等方面是十分有效的,但在系统分析方面有不足之处:

\begin{itemize}
    \item 对时间函数限制严格,绝对可积是充分条件,不少函数不能直接按定义求,如$e^{\alpha t}, \alpha > 0$,其傅里叶变换不存在.

    \item  求解傅里叶反变换比较麻烦.
\end{itemize}
为了解决上述问题拓宽应用范围,人们发现适当地改造满足傅里叶变换的条件.
拉普拉斯变换常用于初始值问题, 即已知某个物理量在初始时刻 $t=0$ 的 值 $f(0)$, 而求解它在初始时刻之后的变化情况 $f(t)$.
 至于它在初始时刻之前 的值, 我们置
$$
f(t)=0, \quad(t<0) .
$$
构造一个带收敛因子$e^{-\sigma t}, \sigma > 0$的函数
$$
g(t) =e^{-\sigma t} f(t)
$$
则$g(t)$绝对可积.于是对$g(t)$进行傅里叶变换:
\begin{equation}
    g(\omega) = \frac{1}{2\pi}  \int_{-\infty}^{\infty} g(t) e^{-\imath \omega t} dt 
    = \frac{1}{2\pi}  \int_{0}^{\infty} f(t) e^{-(\sigma + \imath \omega) t} dt  
\end{equation}
将$\sigma + \imath \omega$记作$p$,令$\bar{f}(p)=2\pi g(\omega)$,则有
函数$f(t)$的拉普拉斯变换$\bar{f}(p)$为
\begin{equation}
    \bar{f}(p) = \mathcal{L} \{ f(t) \} = \int_0 ^{\infty} e^{-pt} f(t) dt .
\end{equation}
该积分称为\textbf{拉普拉斯积分},$\mathcal{L}$称为\textbf{拉普拉斯变换算符}.$e^{-pt}$称为拉普拉斯变换的\textbf{核}.注意
这里的积分上下限.

$g(\omega)$ 的傅里叶逆变换是
$$
g(t)=\int_{-\infty}^{\infty} g(\omega) e^{\imath \omega t} d \omega=\frac{1}{2 \pi} \int_{-\infty}^{\infty} \bar{f}(\sigma+\imath \omega) e^{\imath \omega t} d \omega,
$$
即
$$
f(t)=\frac{1}{2 \pi} \int_{-\infty}^{\infty} \bar{f}(\sigma+\imath \omega) e^{(\sigma+\imath \omega)t} d \omega
$$
由 $\sigma+\imath \omega=p$, 有 $d \omega=\frac{1}{\imath} d p$. 所以
\begin{equation}
    \label{eq:fourier_mellin_integral}
    f(t)=\frac{1}{2 \pi \imath} \int_{\sigma-\imath \infty}^{\sigma+\imath \infty} \bar{f}(p) e^{ p t} d p .
\end{equation}
$\bar{f}(p)$ 又称为\textbf{像函数}, 而 $f(t)$ 称为\textbf{原函数}, 它们之间的关系常用简单的符号写为
\begin{equation}
\begin{aligned}
& \bar{f}(p)=\mathcal{L}[f(t)], \\
& f(t)=\mathcal{L}^{-1}[\bar{f}(p)]
\end{aligned}
\end{equation}
或
\begin{equation}
\begin{aligned}
& \bar{f}(p) \risingdotseq f(t), \\
& f(t) \fallingdotseq \bar{f}(p) .
\end{aligned}
\end{equation}


拉普拉斯变换存在的充分条件是 
\begin{enumerate}
    \item[(1)]  在 $0 \leqslant t<\infty$ 的任一有限区间上; 除了有限个间断点外, 函数 $f(t)$ 及其导数是处处连续的;
    \item[(2)]  存在常数 $M>0$ 和 $\sigma \geqslant 0$, 使对任何 $t$ 值 $(0 \leqslant t<\infty)$, 有
    $$
    |f(t)|<M e^{\sigma t} \text {. }
    $$
\end{enumerate}
$\sigma$ 的下界称为收敛横标, 用 $\sigma_0$ 表示. 在实际应用中, 大多数函数都满足 这个充分条件.
但一个反例是$e^{t^2}$. 对于满足拉普拉斯变换存在条件的函数,通过证明可以知道$\lim_{p\to \infty} \bar{f}(p) = 0$.
也就意味着,若$\bar{f}(p)$在$p\to \infty$的渐近行为是$p$的正幂次,那么逆变换不存在. 且可以证明$\bar{f}(p)$在$\Re \sigma > \sigma_0$的半平面是解析的(具体证明
见梁昆淼书).
不难看出拉普拉斯变换为\textbf{线性变换}, 即满足
$$
\mathcal{L}\{a f(t)+b g(t)\}=a \mathcal{L}\{f(t)\}+b \mathcal{L}\{g(t)\} \,.
$$
不言而喻的是, 傅里叶变换也是线性变换.

对于$\delta$函数,易知
$$
\mathcal{L}\left\{\delta\left(t-t_0\right)\right\}=\int_0^{\infty} e^{-p t} \delta\left(t-t_0\right) d t=e^{-p t_0},  \quad t_0>0 .
$$
% 通过反复迭代
% 可以验证
% $$
% \mathcal{L}\left\{F^{(n)}(t)\right\}=s^n \mathcal{L}\{F(t)\}-s^{n-1} F(+0)-\cdots-F^{(n-1)}(+0) .
% $$
\begin{example}
    求以下函数的拉普拉斯变换
    \begin{enumerate}
        \item $f(t) = 1$,
        \item $f(t) = t$,
        \item $f(t) =e^{s t}, \textrm{s为常数}$,
        \item $f(t) = t e^{s t}, \textrm{s为常数}$,
        \item $f(t) = \cosh {k t}$, $g(t)= \sinh {kt}$.
    \end{enumerate}
\end{example}

\begin{solution}
        


    \begin{enumerate}
        \item $\Re p > 0$, 
            $$\int_{0}^{\infty} 1 \cdot e^{-p t} dt = \frac{1}{p}$$
        \item  $\Re p > 0$, 
            $$\int_{0}^{\infty} t \cdot e^{-p t} dt = \frac{1}{p^2}$$
            因此类似有$\mathcal{L}[t^n] = \frac{n!}{p^{n+1}}$.
        \item $$
            \int_{0}^{\infty} e^{s t} e^{-p t} dt = \frac{1}{p-s} , 
        $$
        要求$\Re (p-s) > 0$.
        \item 
        $$
        \int_{0}^{\infty} t e^{s t} e^{-p t} dt = \frac{1}{(p-s)^2} ,
        $$
        要求$\Re (p-s) > 0$.  
    类似有 $\mathcal{L}[t^n e^{s t}] = \frac{n!}{(p-s)^{n+1}}$.
        \item 利用
        $$
         \cosh{k t}  = \frac{1}{2} \left[ e^{kt} + e^{-kt}\right] \quad 
         \sinh{k t}  = \frac{1}{2} \left[ e^{kt} - e^{-kt}\right] 
        $$

        $$
        \mathcal{L}[ \cosh{k t} ] = \half \left( \frac{1}{p-k} + \frac{1}{p+k} \right) = \frac{p}{p^2 - k^2} ,
        $$
        $$
        \mathcal{L}[ \sinh{k t} ] = \half \left( \frac{1}{p-k} - \frac{1}{p+k} \right) = \frac{k}{p^2 - k^2} .
        $$
    \end{enumerate}
\end{solution}
更多的拉普拉斯变换对见附录二(梁昆淼书).

\subsection{拉普拉斯变换的性质}
与傅里叶变换类似,拉氏变换$\bar{f}(p)$也满足很多基本性质.
\begin{enumerate}
    \item 导数定理
        \begin{equation}
        f^{\prime}(t) \fallingdotseq  p \bar{f}(p)-f(0),
        \end{equation}
        对于高阶导数有递推关系
        \begin{equation}
            f^{(n)}(t) \fallingdotseq p^n \bar{f}(p)-p^{n-1} f(0)-p^{n-2} f^{\prime}(0)-\cdots-p f^{(n-2)}(0)-f^{(n-1)}(0)
        \end{equation}
        \label{eq:lapl_trans_derivative}
    \item 积分定理
        \begin{equation}
        \int_0^t \psi(\tau) d \tau \fallingdotseq  \frac{1}{p} \mathcal{L}[\psi(t)]
        \end{equation}
    \item 相似性定理 
    \begin{equation}
         f(a t) \fallingdotseq  
         \frac{1}{a} \bar{f}\left(\frac{p}{a}\right)
    \end{equation}
        
    \item 位移定理 
        \begin{equation}       
         e^{-\lambda t} f(t) \fallingdotseq  \bar{f}(p+\lambda)
        \end{equation}

    \item 延迟定理 
    \begin{equation}       
        f\left(t-t_0\right) \fallingdotseq  e^{-p t_0} \bar{f}(p).
    \end{equation}

    \item 卷积定理 若  $f_1(t) \fallingdotseq  \bar{f}_1(p), f_2(t) \fallingdotseq  \bar{f}_2(p)$ , 则 
    \begin{equation}       
        f_1(t) * f_2(t) \fallingdotseq  \bar{f}_1(p) \bar{f}_2(p),
    \end{equation}
    其中
    \begin{equation}        
    f_1(t) * f_2(t) \equiv \int_0^t f_1(\tau) f_2(t-\tau) d \tau
    \end{equation}
    称为 $f_1(t)$ 与 $f_2(t)$ 的\textbf{卷积}.
\end{enumerate}

% 若 $f(t)$ 满足 Laplace 变换存在的充分条件, 则
% $$\bar{f}(p) \rightarrow 0, \Re p=s \rightarrow+\infty$$.



% \begin{table}
%     \centering
% \begin{tabular}{|c|c|c|}
%     \hline No & 原 函 数 & 像 函 数 \\
%     \hline 1 & 1 & $\frac{1}{p}$ \\
%     3 & $t^n(n$ 为整数 $)$ & $\frac{n !}{p^{n+1}}$ \\
%     4 & $t^a(a>-1)$ & $\frac{\Gamma(a+1)}{p^{a+1}}$ \\
%     5 & $\sin \omega t$ & $\frac{1}{p-\lambda}$ \\
%     6 & $\cos \omega t$ & $\frac{\omega}{p^2+\omega^2}$ \\
%     7 & $\operatorname{sh} \omega t$ & $\frac{p}{p^2+\omega^2}$ \\
%     8 & $\operatorname{ch} \omega t$ & $\frac{\omega}{p^2-\omega^2}$ \\
%     & & $\frac{p}{p^2-\omega^2}$ \\
%     9 & $\frac{\omega}{(p+\lambda)^2+\omega^2}$ \\
%     10 & $e^{-\lambda t} \sin \omega t$ & $\frac{p+\lambda}{(p+\lambda)^2+\omega^2}$ \\
%     11 & $e^{-\lambda t} t^a$ & $\frac{\Gamma(a+1)}{(p+\lambda)^{a+1}}$ \\
%     12 & $\frac{1}{\sqrt{\pi t}}$ & $\frac{1}{\sqrt{p}}$ \\
%     13 & $\frac{1}{\sqrt{\pi t}} e^{-a^{2 / 4 t}}$ & $\frac{e^{-a \sqrt{p}}}{\sqrt{p}}$ \\
%     14 & $\frac{1}{\sqrt{\pi t}} e^{-2 a \sqrt{t}}$ & $\frac{1}{\sqrt{p}} e^{\frac{a^2}{p}} erfc\left(\frac{a}{\sqrt{p}}\right)$ \\
%     15 & $\frac{1}{\sqrt{\pi a}} \sin 2 \sqrt{a t}$ & $\frac{1}{p \sqrt{p}} e^{-\frac{a}{p}}$ \\
%     16 & $\frac{1}{\sqrt{\pi a}} \cos 2 \sqrt{a t}$ \\
%     17 & $erf(\sqrt{a t})$ & $\frac{1}{\sqrt{p}} e^{-\frac{a}{p}}$ \\
%     18 & $erfc\left(\frac{a}{2 \sqrt{t}}\right)$ & $\frac{\sqrt{a}}{p \sqrt{p+a}}$ \\
%     19 & $e^t erfc(\sqrt{t})$ & $\frac{1}{p} e^{-a \sqrt{p}}$ \\
%     20 & $\frac{1}{\sqrt{\pi t}}-e^t erfc(\sqrt{t})$ & $\frac{1}{p+\sqrt{p}}$ \\
%     21 & $\frac{1}{\sqrt{\pi t}} e^{-a t}+\sqrt{a} \operatorname{erf} \sqrt{a t}$ & $\frac{1}{1+\sqrt{p}}$ \\
%     & $\frac{\sqrt{p+a}}{p}$ \\
%     22 & $J_0(t)$ & $\frac{1}{\sqrt{p^2+1}}$ \\
%     23 & $J_n(t)$ & $\frac{\left(\sqrt{p^2+1}-p\right)^n}{\sqrt{p^2+1}}$ \\
%     24 & $\frac{J_n(a t)}{t}$ & $\frac{1}{n a^n}\left(\sqrt{p^2+a^2}-p\right)^n$ \\
%     26 & $e^{-a t} I_0(b t)$ & $\frac{1}{\sqrt{(p+a)^2-b^2}}$ \\
%     27 & $\lambda^n e^{-\lambda t} I_n(\lambda t)$ & $\frac{\left\{\sqrt{p^2+2 \lambda p}-(p+\lambda)\right\}^n}{\sqrt{p^2+2 \lambda p}}$ \\
%     28 & $t^n J_n(t)\left(n>-\frac{1}{2}\right)$ & $\frac{2^n \Gamma\left(n+\frac{1}{2}\right)}{\sqrt{\pi}} \frac{1}{\left(p^2+1\right)^{n+\frac{1}{2}}}$ \\
%     29 & $t^{\frac{n}{2}} J_n(2 \sqrt{t})$ & $\frac{1}{p^{n+1}} e^{-\frac{1}{p}}$ \\
%     30 & $J_0\left(a \sqrt{t^2-\tau^2}\right) H(t-\tau)$ & $\frac{1}{\sqrt{p^2+a^2}} e^{-\tau \sqrt{p^2+a^2}}$ \\
%     31 & $\frac{J_1\left(a \sqrt{t^2-\tau^2}\right)}{\sqrt{t^2-\tau^2}} H(t-\tau)$ & $\frac{e^{-\tau p}-e^{-\tau \sqrt{p^2+a^2}}}{a \tau}$ \\
%     32 & $\int_t^{\infty} \frac{J_0(t)}{t} d t$ & $\frac{1}{p} \ln \left(p+\sqrt{1+p^2}\right)$ \\
%     33 & $\frac{e^{b t}-e^{a t}}{t}$ & $\ln \frac{p-a}{p-b}$ \\
%     34 & $\frac{1}{\sqrt{\pi t}} \sin \frac{1}{2 t}$ & $\frac{1}{\sqrt{p}} e^{-\sqrt{p}} \sin \sqrt{p}$ \\
%     35 & $\frac{1}{\sqrt{\pi t}} \cos \frac{1}{2 t}$ & $\frac{1}{\sqrt{p}} e^{-\sqrt{p}} \cos \sqrt{p}$ \\
%     \hline
% \end{tabular}
% \end{table}

\subsection{拉普拉斯反变换}
\label{subsec:inverse_laplace_transform}
前面对于拉普拉斯反变换(亦称逆变换)
\begin{equation}
    f(t) = \mathcal{L}^{-1} [ \bar{f}(p)],
\end{equation}
其具体表达式为
\begin{equation}
    f(t)=\frac{1}{2 \pi \imath} \int_{\sigma-\imath \infty}^{\sigma+\imath \infty} \bar{f}(p) e^{ p t} dp, 
\end{equation}
见式\eqref{eq:fourier_mellin_integral}.该式是著名的Bromwich积分(也叫做Riemann-Mellin积分).
由于像函数$\bar{f}(p)$是解析函数是$p$的解析函数,上述积分可以通过留数定理求得.这里需要用到推广的
约旦定理(具体证明略去).

\textbf{推广的约当引理} 设 $C_R$ 是以 $p=0$ 为圆心, 以 $R$ 为半径的圆周在直线 $\Re p=a(>0)$ 左侧的圆弧. 若当 $|p| \rightarrow \infty$ 时, $\bar{f}(p)$ 在 $\frac{\pi}{2}-\delta \leqslant \operatorname{Arg} p \leqslant \frac{3}{2} \pi+\delta$ 中一致趋于零 ( $\delta$ 是小于 $\frac{\pi}{2}$ 的任意正数), 则
$$
\lim _{R \rightarrow \infty} \int_{C_R} \bar{f}(p) e^{p t} d p=0 \quad(t>0) .
$$
由于所作路径是在$\Re p = a$左半平面,而右半平面$\bar{f}(p)$解析, 最终我们得到
\begin{equation}
    f(t) = \sum_{p_j \in \text{全平面}} Res [ \bar{f}(p_j) e^{p_j t}].
\end{equation}

\begin{example}
利用Bromwich积分求解$1/\sqrt{p}$的拉普拉斯变换原函数.
\end{example}
\begin{solution}
    求解参考梁昆淼书.
    被积函数有唯一的奇点 $p=0$, 这是极点型支点, 在考虑回路积分时, 
    不能像简单地加上圆弧 $C_R$, 因为在函数的多值区域不能作积分运算及应用留数定理. 
    为此, 先自 $p=0$ 沿负实轴至 $\infty$ 将 $p$ 平面切割开,
    再作如图回路, 然后在被积函数的一个单值分支上作积分运算,并应用留数定理. 
    被积函数在回路所围区域上无奇点,
    % 按柯西定理
    % 其中 $l_1$ 和 $l_2$ 分别沿着割线的下沿和上沿,
     圆弧 $C_R$ 和圆 $C_{\varepsilon}$均以原点为圆心, 
     半径分别为 $R$ 和 $\varepsilon$. 令 $R \rightarrow \infty$, 
     由推广的约当引理知, $\lim _{R \rightarrow \infty} \int_{C_R}=0$, 
     同时不难验证 $\lim _{s \rightarrow 0} \int_{C_s}=0$.
     于是,
$$
f(t)=\int_{a-\imath \infty}^{a+\imath \infty} \frac{e^{p t}}{\sqrt{p}} d p=-\frac{1}{2 \pi \imath} \lim _{R \rightarrow \infty}\left(\int_{l_1}+\int_{l_2}\right) \frac{e^{p t}}{\sqrt{p}} d p .
$$
$p$ 在割线下沿 $l_1$ 和上沿 $l_2$ 上的辐角分别为 $-\pi$ 和 $\pi$. 这样, 在 $l_1$ 上, $\sqrt{p}=\sqrt{|p|} e^{-\imath \pi / 2}=$ $-\imath \sqrt{|p|}$; 在 $l_2$ 上 $\sqrt{p}=\sqrt{|p|} e^{\imath \pi / 2}=\imath \sqrt{|p|}$. 于是,
$$
\begin{aligned}
f(t) & =-\frac{1}{2 \pi \imath} \int_0^{-\infty} e^{p t} \frac{1}{-\imath \sqrt{|p|}} d p-\frac{1}{2 \pi \imath} \int_{-\infty}^0 e^{p t} \frac{1}{\imath \sqrt{|p|}} d p \\
& =\frac{1}{\pi} \int_{-\infty}^0 e^{p t} \frac{1}{\sqrt{|p|}} d p .
\end{aligned}
$$
改用 $y=\sqrt{|p| t}$ 作为积分变数,
$$
\begin{aligned}
f(t) & =\frac{1}{\pi} \int_{\infty}^0 e^{-y^2} \frac{\sqrt{t}}{y}\left(-\frac{2 y}{t} d y\right) \\
& =\frac{2}{\pi \sqrt{t}} \int_0^{\infty} e^{-y^2} d y=\frac{2}{\pi \sqrt{t}} \cdot \frac{\sqrt{\pi}}{2}=\frac{1}{\sqrt{\pi t}} .
\end{aligned}
$$
\end{solution}

很多时候,我们并不总需要用上述方法求解,对于有些问题,可以使用利用拉普拉斯变换的性质结合常见函数的表达式来求解.
对于像函数的导数, $f(t)$是满足Laplace变换的充分条件,那么$\bar{f}(p)$在
$\Re p > \sigma_0$的半平面解析,
于是求$n$阶导数得
\begin{equation}
\bar{f}^{(n)}(p)=\frac{d^{n}}{d p^n} \int_0^{\infty} f(t) e^{-p t} d t=\int_0^{\infty}(-t)^n f(t) e^{-p t} d t
\end{equation}
所以有
\begin{equation}
    \bar{f}^{(n)}(p)\risingdotseq (-t)^n f(t)
\end{equation}
不难发现,

$$    \begin{aligned}
    & \frac{1}{p^2}=-\frac{d}{d p} \frac{1}{p} \risingdotseq  t \\
    & \frac{1}{p^3}=\frac{1}{2} \frac{d^2}{d p^2} \frac{1}{p} \risingdotseq \frac{1}{2} t^2 .
    \end{aligned}$$


例如, 利用留数性质可以用待定系数法分解
$$
    \begin{aligned}
    \frac{1}{p^3(p+\alpha)} & =\frac{1}{\alpha} \frac{1}{p^3}-\frac{1}{\alpha^2} \frac{1}{p^2}+\frac{1}{\alpha^3} \frac{1}{p}-\frac{1}{\alpha^3} \frac{1}{p+\alpha} \\
    & \risingdotseq \frac{1}{2 \alpha} t^2-\frac{1}{\alpha^2} t+\frac{1}{\alpha^3}-\frac{1}{\alpha^3} e^{-\alpha t}
    \end{aligned}
$$

\begin{example}
求 $\frac{e^{-\alpha p}}{p(p+b)}$ 的原函数.
\end{example}
\begin{solution}
    $\frac{1}{p} \risingdotseq \Theta(t)$, 应用延迟定理,
    $$
e^{-\alpha p} / p \risingdotseq \Theta(t-\alpha) \text {. 又 } 1 /(p+b) \risingdotseq e^{-b t} \text {. }
$$

这样, 本题给的像函数可以看作 $e^{-\alpha p} / p$ 与 $1 /(p+b)$ 的乘积, 应用卷积定理,即得
$$
\begin{aligned}
\frac{e^{-\alpha p}}{p(p+b)} & \risingdotseq \int_0^t \Theta(\tau-\alpha) e^{-b(t-\tau)} d \tau \\
& =\Theta(t-\alpha) \int_\alpha^t e^{-b(t-\tau)} d \tau \\
& =\Theta(t-\alpha)\left[\frac{1}{b} e^{-b(t-\tau)}\right]_\alpha^t \\
& =\frac{1}{b}\left[1-e^{-b(t-\alpha)}\right] \Theta(t-\alpha) .
\end{aligned}
$$
读者也可以运用裂项分解分母来求解验证.
\end{solution}



\section{积分变换的应用}
\label{sec:integral_transform_applications}
\subsection{拉普拉斯积分变换的应用}
\label{subsec:applications}
拉普拉斯变换有很多应用,这里从三个方面的例子来说明.
\subsubsection{计算级数和}
有时,拉普拉斯变换可以用来计算某些级数的和.以下面的例子说明.

如计算级数和
$$
    \sum_{n=1} ^{\infty}  \frac{1}{n^2},
$$
由前面例题
$$
\int_0^{\infty} t e^{-p t} d t=\frac{1}{p^2}, \quad \Re p>0
$$

将级数化为
$$
\begin{aligned}
\sum_{n= 1}^{\infty} \frac{1}{n^2} & =\sum_{n=1}^{\infty} \int_0^{\infty} t e^{-n t} d t \\
& =\int_0^{\infty} t\left[\sum_{n=1}^{\infty} e^{-n t}\right] d t=\int_0^{\infty} \frac{t}{e^t-1} d t
\end{aligned}
$$
这个积分的值是 $\pi^2/6$,下面用 Parseval 等式核对,不必查表.

在 $[-\pi,\pi]$ 上取 $f(x)=x$,则 $\ell=\pi$,
\begin{equation}
  c_n = \frac{1}{2\pi}\int_{-\pi}^{\pi} x\, e^{-\imath n x}\,\mathrm{d}x.
\end{equation}
$n\neq 0$ 时分部积分.$\int_{-\pi}^{\pi} e^{-\imath n x}\,\mathrm{d}x=0$,只剩边界项.两端 $e^{-\imath n\pi}=e^{\imath n\pi}=(-1)^n$,并且 $1/(-\imath)=\imath$,
\begin{equation}
  \left[\frac{x\,e^{-\imath n x}}{-\imath n}\right]_{-\pi}^{\pi}
  = \frac{2\pi\,\imath}{n}(-1)^n.
\end{equation}
再除以 $2\pi$,得到 $c_n = \imath (-1)^n/n$,$|c_n|^2=1/n^2$,而 $c_0=0$.Parseval 等式
\begin{equation}
  \sum_{n=-\infty}^{\infty}|c_n|^2
  = \frac{1}{2\pi}\int_{-\pi}^{\pi} x^2\,\mathrm{d}x.
\end{equation}
左边是 $2\sum_{n=1}^{\infty} 1/n^2$.右边是
\begin{equation}
  \frac{1}{2\pi}\cdot\frac{2\pi^3}{3} = \frac{\pi^2}{3}.
\end{equation}
因此 $\sum_{n=1}^{\infty} 1/n^2=\pi^2/6$.左边的积分是 $\pi^2/3$,再除以 $2$ 才得到 $\pi^2/6$.

\subsubsection{求解定积分}
如果 $\int_v^{\infty} \bar{f}(q) d q$ 存在, 且当$t\to 0$时, $|f(t)/t|$有界,则
\begin{equation}
    \int_{p}^{\infty} \bar{f}(q) dq \risingdotseq \frac{f(t)}{t} . 
\end{equation}
比如
$$
\frac{\sin \omega t}{t} \fallingdotseq \int_p^{\infty} \frac{\omega}{q^2+\omega^2} d q=\frac{\pi}{2}-\arctan \frac{p}{\omega}.
$$

当$p\to 0$时, 有
\begin{equation}
    \int_0^{\infty} \bar{f}(p) d p=\int_0^{\infty} \frac{f(t)}{t} d t
\end{equation}
一个例子如
$$
\int_0^{\infty} \frac{\sin t}{t} d t=\int_0^{\infty} \frac{1}{p^2+1} d p=\frac{\pi}{2}
$$
不仅如此, 有些积分无法用留数定理计算,如
$$
\int_0^{\infty} \frac{\cos a t-\cos b t}{t} d t \quad a>0, b>0
$$
$\mathcal{L}\{\cos at\}=p/(p^2+a^2)$.于是
\begin{equation}
  \int_0^{\infty}\frac{\cos(at)-\cos(bt)}{t}\,\mathrm{d}t
  = \int_0^{\infty}\left(\frac{p}{p^2+a^2}-\frac{p}{p^2+b^2}\right)\mathrm{d}p.
\end{equation}
不定积分是 $\half\ln(p^2+a^2)$.定积分
\begin{equation}
  \left[\half\ln\frac{p^2+a^2}{p^2+b^2}\right]_0^{\infty}
  = 0 - \half\ln\frac{a^2}{b^2}
  = \ln\frac{b}{a}.
\end{equation}
上限处分子分母都像 $p^2$,对数趋于 $\ln 1=0$.

\subsubsection{求解微分方程}
给一个例子,利用拉普拉斯变换求解简谐振子方程的解.
\begin{example}
利用拉普拉斯变换求解简谐振子方程的解.
\end{example}
\begin{solution}
    质量为$m$的质点在弹性系数为$k$的弹簧牵引下做简谐运动,满足方程为
    $$
m \frac{d^2 X(t)}{d t^2}+k X(t)=0
$$
初始条件取
$$
X(0)=X_0, \quad X^{\prime}(0)=0
$$
应用拉氏变换到该方程上得到
$$
m \mathcal{L}\left\{\frac{d^2 X}{d t^2}\right\}+k \mathcal{L}\{X(t)\}=0
$$
用$x(p)$表示未知变换$\mathcal{L}\{ X(t)\}$,于是根据导数定理\eqref{eq:lapl_trans_derivative}有
$$
m p^2 x(p)-m p X_0+k x(p)=0,
$$
化简为
$$
x(p)=X_0 \frac{p}{p^2+\omega_0^2}, \quad \omega_0^2 \equiv \frac{k}{m} .
$$
查表得
$$
X(t)=X_0 \cos \omega_0 t.
$$
初速度为零,解里没有正弦项.
\end{solution}

若初速度改为 $X'(0)=V_0$,二阶导数定理多一项 $-m V_0$,于是
\begin{equation}
    x(p)=X_0\frac{p}{p^2+\omega_0^2}+V_0\frac{1}{p^2+\omega_0^2},
\end{equation}
\begin{equation}
    X(t)=X_0\cos(\omega_0 t)+\frac{V_0}{\omega_0}\sin(\omega_0 t).
\end{equation}
正弦项的系数是 $V_0/\omega_0$,因为 $\mathcal{L}\{\sin(\omega_0 t)\}=\omega_0/(p^2+\omega_0^2)$.在 $t=0$ 对解求导,可以核对这个因子.
傅里叶级数的系数要把积分算完.傅里叶变换按正文取极限后的约定
\begin{equation*}
  F(\omega)=\frac{1}{2\pi}\int_{-\infty}^{\infty}f(x)e^{-\mathrm{i}\omega x}\,\mathrm{d}x,
  \qquad
  f(x)=\int_{-\infty}^{\infty}F(\omega)e^{\mathrm{i}\omega x}\,\mathrm{d}\omega.
\end{equation*}
拉普拉斯变换用 $\bar f(p)=\int_0^{\infty}e^{-pt}f(t)\,\mathrm{d}t$.

\begin{problemset}
\renewcommand{\theenumi}{\thechapter.\arabic{enumi}}
\renewcommand{\labelenumi}{\theenumi.}
\item 对正整数 $m,n\geq 1$,在 $[0,\ell]$ 上验证
\begin{equation*}
  \int_0^{\ell}
  \cos\Bigl(\frac{n\pi x}{\ell}\Bigr)
  \cos\Bigl(\frac{m\pi x}{\ell}\Bigr)
  \,\mathrm{d}x
  =\frac{\ell}{2}\,\delta_{mn},
\end{equation*}
并单独计算 $n=0$ 时 $\int_0^{\ell}1\,\mathrm{d}x$.由此说明半幅余弦级数里 $a_0$ 和 $a_n$ 可以共用因子 $2/\ell$,常数项仍然写成 $a_0/2$.

\item 热棒两端保持零度,初始温度 $f(x)=x(\ell-x)$.求出正弦系数,并写出 $T(x,t)$.

\item 正文里 $f(x)=x^2$ 的 Parseval 等式左边是 $\pi^4/5$,右边含 $\sum 1/n^4$.两边各自算完,得到 $\sum_{n=1}^{\infty}1/n^4=\pi^4/90$.

\item 在 $[-\pi,\pi]$ 上把 $f(x)=x$ 展成傅里叶级数.按正文的能量定理,分别计算 $\frac{1}{2\pi}\int_{-\pi}^{\pi}x^2\,\mathrm{d}x$ 和 $\sum b_n^2/2$,从而得到 $\sum_{n=1}^{\infty}1/n^2=\pi^2/6$.两边都要单独算完再比较.

\item 按上面的傅里叶变换约定计算 $\mathcal{F}[e^{-\alpha|x|}]$, $\alpha>0$.再用反演积分在 $x=0$ 处核对 $f(0)=1$.

\item 计算宽度为 $2a$、中心在原点的矩形脉冲的变换,再计算中心移到 $x_0$ 的同一脉冲.哪些 $\omega$ 使中心在原点的变换为零?平移为什么不改变这些零点?

\item 用分部积分证明 $\mathcal{F}[f']=\mathrm{i}\omega F$,并写明边界项为什么必须为零.对全直线上的 $u_t=u_{xx}$, $u$ 在无穷远处趋于零,写出 $\tilde u(\omega,t)$ 满足的方程.

\item 由 $\mathcal{L}\{e^{\mathrm{i}\omega t}\}$ 得到 $\mathcal{L}\{\sin(\omega t)\}$ 和 $\mathcal{L}\{\cos(\omega t)\}$,并写明收敛的半平面.再用位移定理计算 $\mathcal{L}\{e^{-\lambda t}\sin(\omega t)\}$,不要重新积分.

\item 用拉普拉斯变换解 $X''-3X'+2X=e^{t}$, $X(0)=1$, $X'(0)=0$.部分分式反演后,在结果上核对两个初值.

\item 反演 $e^{-2p}/[(p+1)(p+2)]$.分别用卷积和用部分分式加延迟定理,并确认 $0\leq t<2$ 时原函数是零.
\end{problemset}
